\textbf{File System}
\begin{lstlisting}[style=CodeListing]
#ifndef UNICODE
#define UNICODE
#endif
#ifndef _UNICODE
#define _UNICODE
#endif

#include "../../include/filesystem.h"
#include "../../include/utils.h"
#include <windows.h>
#include <commctrl.h>
#include <shellapi.h>
#include <algorithm>

namespace Filesystem {
    bool HasSubdirs(const std::wstring &path) {
        std::wstring pattern = Utils::JoinPath(path, L"*");
        WIN32_FIND_DATAW fd{};
        HANDLE h = FindFirstFileW(pattern.c_str(), &fd);
        if (h == INVALID_HANDLE_VALUE) return false;
        bool result = false;
        do {
            if ((fd.dwFileAttributes & FILE_ATTRIBUTE_DIRECTORY) == 0) continue;
            if (wcscmp(fd.cFileName, L".") == 0 || wcscmp(fd.cFileName, L"..") == 0) continue;
            result = true;
            break;
        } while (FindNextFileW(h, &fd));
        FindClose(h);
        return result;
    }
    
    void PopulateListForDirectory(HWND hList, const std::wstring &dirPath, 
                                  std::vector<ListItem> &items, HWND hEmptyLabel) {
        ListView_DeleteAllItems(hList);
        items.clear();
        
        // Normalize path - resolve any symbolic links or short paths
        std::wstring normalizedPath = dirPath;
        wchar_t longPath[MAX_PATH + 1] = {0};
        DWORD longPathLen = GetLongPathNameW(dirPath.c_str(), longPath, MAX_PATH);
        if (longPathLen > 0 && longPathLen <= MAX_PATH) {
            normalizedPath = longPath;
        }
        
        // Ensure path ends with backslash for directories (except root like C:\)
        if (normalizedPath.length() > 3 && normalizedPath.back() != L'\\') {
            normalizedPath += L"\\";
        }
        
        // Check if path exists - try original path first, then normalized
        DWORD attribs = GetFileAttributesW(normalizedPath.c_str());
        if (attribs == INVALID_FILE_ATTRIBUTES && normalizedPath != dirPath) {
            // Try original path if normalized failed
            attribs = GetFileAttributesW(dirPath.c_str());
            if (attribs != INVALID_FILE_ATTRIBUTES) {
                normalizedPath = dirPath;
            }
        }
        
        if (attribs == INVALID_FILE_ATTRIBUTES) {
            // Path doesn't exist - check if it's a network path that might need time to connect
            if (hEmptyLabel) {
                // For network paths, show "Подключение..." instead of "Файл не найден"
                bool isNetworkPath = (normalizedPath.length() >= 2 && normalizedPath[0] == L'\\' && normalizedPath[1] == L'\\');
                if (isNetworkPath) {
                    SetWindowTextW(hEmptyLabel, L"Подключение...");
                } else {
                    SetWindowTextW(hEmptyLabel, L"Файл не найден");
                }
                ShowWindow(hEmptyLabel, SW_SHOW);
                SetWindowPos(hEmptyLabel, HWND_TOP, 0, 0, 0, 0, SWP_NOMOVE | SWP_NOSIZE);
                InvalidateRect(hEmptyLabel, nullptr, TRUE);
            }
            // Hide list view to show only the message
            ShowWindow(hList, SW_HIDE);
            return;
        }
        
        // Check if it's a directory
        if (!(attribs & FILE_ATTRIBUTE_DIRECTORY)) {
            // It's a file, not a directory - show "Файл пуст" message
            if (hEmptyLabel) {
                SetWindowTextW(hEmptyLabel, L"Файл пуст");
                ShowWindow(hEmptyLabel, SW_SHOW);
                SetWindowPos(hEmptyLabel, HWND_TOP, 0, 0, 0, 0, SWP_NOMOVE | SWP_NOSIZE);
                InvalidateRect(hEmptyLabel, nullptr, TRUE);
            }
            // Hide list view to show only the message
            ShowWindow(hList, SW_HIDE);
            InvalidateRect(hList, nullptr, TRUE);
            UpdateWindow(hList);
            RedrawWindow(hList, nullptr, nullptr, RDW_INVALIDATE | RDW_UPDATENOW | RDW_ERASE);
            return;
        }

        // Check if it's a network path - these may need time to load
        bool isNetworkPath = (normalizedPath.length() >= 2 && normalizedPath[0] == L'\\' && normalizedPath[1] == L'\\');
        
        std::wstring pattern = Utils::JoinPath(normalizedPath, L"*");
        WIN32_FIND_DATAW fd{};
        HANDLE h = FindFirstFileW(pattern.c_str(), &fd);
        
        // If first attempt fails, show loading and retry after a delay
        // This handles cases where directory is still loading (especially network paths and special folders)
        if (h == INVALID_HANDLE_VALUE) {
            DWORD error = GetLastError();
            
            // Show loading message while we retry
            if (hEmptyLabel) {
                SetWindowTextW(hEmptyLabel, L"Загрузка...");
                ShowWindow(hEmptyLabel, SW_SHOW);
                SetWindowPos(hEmptyLabel, HWND_TOP, 0, 0, 0, 0, SWP_NOMOVE | SWP_NOSIZE);
                ShowWindow(hList, SW_HIDE);
                InvalidateRect(hEmptyLabel, nullptr, TRUE);
                UpdateWindow(hList);
            }
            
            // Retry after delay - network paths need more time, local paths also need time for special folders
            // Special folders like Downloads might need time to initialize
            int delay = isNetworkPath ? 1500 : 800; // 1.5 seconds for network, 800ms for local (handles special folders)
            
            // Use a simple delay that doesn't block too much
            // But we need to process messages to keep UI responsive
            DWORD startTime = GetTickCount();
            while ((GetTickCount() - startTime) < delay) {
                // Process a few messages to keep UI responsive
                MSG msg;
                while (PeekMessageW(&msg, nullptr, 0, 0, PM_REMOVE)) {
                    TranslateMessage(&msg);
                    DispatchMessageW(&msg);
                    if ((GetTickCount() - startTime) >= delay) break;
                }
                if ((GetTickCount() - startTime) < delay) {
                    Sleep(50); // Small sleep to avoid busy waiting
                }
            }
            
            // Try again
            h = FindFirstFileW(pattern.c_str(), &fd);
        }
        
        if (h == INVALID_HANDLE_VALUE) {
            // Directory is inaccessible or empty after retry
            DWORD error = GetLastError();
            if (hEmptyLabel) {
                if (error == ERROR_ACCESS_DENIED || error == ERROR_SHARING_VIOLATION) {
                    SetWindowTextW(hEmptyLabel, L"Доступ запрещён");
                } else if (isNetworkPath) {
                    // For network paths, even after retry, show loading - they might still be connecting
                    SetWindowTextW(hEmptyLabel, L"Загрузка...");
                } else if (error == ERROR_FILE_NOT_FOUND || error == ERROR_NO_MORE_FILES) {
                    // For local paths after retry, if still empty, show empty message
                    SetWindowTextW(hEmptyLabel, L"Здесь пусто");
                } else {
                    // Other local errors - assume empty
                    SetWindowTextW(hEmptyLabel, L"Здесь пусто");
                }
                ShowWindow(hEmptyLabel, SW_SHOW);
                SetWindowPos(hEmptyLabel, HWND_TOP, 0, 0, 0, 0, SWP_NOMOVE | SWP_NOSIZE);
                InvalidateRect(hEmptyLabel, nullptr, TRUE);
            }
            // Hide list view to show only the message
            ShowWindow(hList, SW_HIDE);
            InvalidateRect(hList, nullptr, TRUE);
            UpdateWindow(hList);
            RedrawWindow(hList, nullptr, nullptr, RDW_INVALIDATE | RDW_UPDATENOW | RDW_ERASE);
            return;
        }
        
        int idx = 0;
        do {
            std::wstring name = fd.cFileName;
            if (name == L"." || name == L"..") continue;
            bool isDir = (fd.dwFileAttributes & FILE_ATTRIBUTE_DIRECTORY) != 0;

            ListItem ri{ name, Utils::JoinPath(normalizedPath, name), isDir };
            items.push_back(ri);

            LVITEMW item{};
            SHFILEINFOW sfi{};
            UINT attrs = isDir ? FILE_ATTRIBUTE_DIRECTORY : FILE_ATTRIBUTE_NORMAL;
            SHGetFileInfoW(ri.fullPath.c_str(), attrs, &sfi, sizeof(sfi), 
                          SHGFI_SYSICONINDEX | SHGFI_USEFILEATTRIBUTES | SHGFI_SMALLICON);

            item.mask = LVIF_TEXT | LVIF_PARAM | LVIF_IMAGE;
            item.iItem = idx;
            item.pszText = const_cast<wchar_t*>(name.c_str());
            item.lParam = static_cast<LPARAM>(idx);
            item.iImage = sfi.iIcon;
            ListView_InsertItem(hList, &item);

            std::wstring typeText = isDir ? L"Dir" : L"File";
            ListView_SetItemText(hList, idx, 1, const_cast<wchar_t*>(typeText.c_str()));

            std::wstring mod = Utils::FileTimeToWString(fd.ftLastWriteTime);
            ListView_SetItemText(hList, idx, 2, const_cast<wchar_t*>(mod.c_str()));
            ++idx;
        } while (FindNextFileW(h, &fd));
        FindClose(h);
        
        if (hEmptyLabel) {
            if (items.empty()) {
                // Directory is empty - show message and hide list view
                SetWindowTextW(hEmptyLabel, L"Здесь пусто");
                ShowWindow(hEmptyLabel, SW_SHOW);
                SetWindowPos(hEmptyLabel, HWND_TOP, 0, 0, 0, 0, SWP_NOMOVE | SWP_NOSIZE);
                InvalidateRect(hEmptyLabel, nullptr, TRUE);
                // Hide list view to show only the message
                ShowWindow(hList, SW_HIDE);
            } else {
                // Directory has items - show list view and hide message
                ShowWindow(hEmptyLabel, SW_HIDE);
                ShowWindow(hList, SW_SHOW);
            }
        } else {
            // If no empty label, always show list view
            ShowWindow(hList, SW_SHOW);
        }
        
        InvalidateRect(hList, nullptr, TRUE);
        UpdateWindow(hList);
        RedrawWindow(hList, nullptr, nullptr, RDW_INVALIDATE | RDW_UPDATENOW | RDW_ERASE);
    }
    
    HTREEITEM InsertTreeItem(HWND hTree, HTREEITEM hParent, 
                             const std::wstring &text, const std::wstring &fullPath, 
                             bool placeholderChild) {
        TVINSERTSTRUCTW ins{};
        ins.hParent = hParent;
        ins.hInsertAfter = TVI_LAST;
        ins.item.mask = TVIF_TEXT | TVIF_PARAM;
        ins.item.pszText = const_cast<wchar_t*>(text.c_str());
        NodeData *data = new NodeData{fullPath, false};
        ins.item.lParam = reinterpret_cast<LPARAM>(data);
        HTREEITEM hItem = TreeView_InsertItem(hTree, &ins);
        if (placeholderChild) {
            TVINSERTSTRUCTW child{};
            child.hParent = hItem;
            child.hInsertAfter = TVI_LAST;
            child.item.mask = TVIF_TEXT;
            child.item.pszText = const_cast<wchar_t*>(L"...");
            TreeView_InsertItem(hTree, &child);
        }
        return hItem;
    }
    
    void ExpandTreeItem(HWND hTree, HTREEITEM hItem) {
        // Remove placeholder children
        HTREEITEM hChild = TreeView_GetChild(hTree, hItem);
        while (hChild) {
            TVITEMW tvi{}; tvi.mask = TVIF_TEXT; tvi.hItem = hChild;
            wchar_t buf[8]; tvi.pszText = buf; tvi.cchTextMax = 8;
            TreeView_GetItem(hTree, &tvi);
            if (wcscmp(buf, L"...") == 0) {
                TreeView_DeleteItem(hTree, hChild);
                break;
            }
            hChild = TreeView_GetNextSibling(hTree, hChild);
        }

        TVITEMW tvi{}; tvi.mask = TVIF_PARAM; tvi.hItem = hItem; TreeView_GetItem(hTree, &tvi);
        NodeData *data = reinterpret_cast<NodeData*>(tvi.lParam);
        if (!data || data->expanded) return;

        std::wstring pattern = Utils::JoinPath(data->path, L"*");
        WIN32_FIND_DATAW fd{};
        HANDLE h = FindFirstFileW(pattern.c_str(), &fd);
        if (h == INVALID_HANDLE_VALUE) return;
        do {
            if ((fd.dwFileAttributes & FILE_ATTRIBUTE_DIRECTORY) == 0) continue;
            std::wstring name = fd.cFileName;
            if (name == L"." || name == L"..") continue;
            std::wstring full = Utils::JoinPath(data->path, name);
            bool hasChildren = HasSubdirs(full);
            InsertTreeItem(hTree, hItem, name, full, hasChildren);
        } while (FindNextFileW(h, &fd));
        FindClose(h);
        data->expanded = true;
    }
    
    void FreeNodeDataRecursive(HWND hTree, HTREEITEM hItem) {
        while (hItem) {
            HTREEITEM child = TreeView_GetChild(hTree, hItem);
            if (child) FreeNodeDataRecursive(hTree, child);
            TVITEMW tvi{}; tvi.mask = TVIF_PARAM; tvi.hItem = hItem; TreeView_GetItem(hTree, &tvi);
            if (tvi.lParam) delete reinterpret_cast<NodeData*>(tvi.lParam);
            hItem = TreeView_GetNextSibling(hTree, hItem);
        }
    }
    
    HTREEITEM EnsurePathInTree(HWND hTree, const std::wstring &fullPath) {
        if (fullPath.size() < 3 || fullPath[1] != L':') return nullptr;
        
        HTREEITEM hItem = TreeView_GetRoot(hTree);
        HTREEITEM driveItem = nullptr;
        while (hItem) {
            TVITEMW tvi{}; tvi.mask = TVIF_PARAM; tvi.hItem = hItem; TreeView_GetItem(hTree, &tvi);
            NodeData *data = reinterpret_cast<NodeData*>(tvi.lParam);
            if (data && _wcsicmp(data->path.c_str(), std::wstring(fullPath.substr(0,3)).c_str()) == 0) {
                driveItem = hItem; break;
            }
            hItem = TreeView_GetNextSibling(hTree, hItem);
        }
        if (!driveItem) return nullptr;

        if (fullPath.size() == 3) return driveItem;

        size_t pos = 3;
        HTREEITEM current = driveItem;
        while (pos < fullPath.size()) {
            size_t next = fullPath.find(L'\\', pos);
            std::wstring part = fullPath.substr(pos, next == std::wstring::npos ? std::wstring::npos : next - pos);
            if (part.empty()) break;

            ExpandTreeItem(hTree, current);

            HTREEITEM child = TreeView_GetChild(hTree, current);
            HTREEITEM foundChild = nullptr;
            while (child) {
                TVITEMW tv{}; tv.mask = TVIF_PARAM; tv.hItem = child; TreeView_GetItem(hTree, &tv);
                NodeData *cd = reinterpret_cast<NodeData*>(tv.lParam);
                if (cd) {
                    const wchar_t *lastSep = wcsrchr(cd->path.c_str(), L'\\');
                    std::wstring last = lastSep ? std::wstring(lastSep + 1) : cd->path;
                    if (_wcsicmp(last.c_str(), part.c_str()) == 0) { foundChild = child; break; }
                }
                child = TreeView_GetNextSibling(hTree, child);
            }
            if (!foundChild) return current;
            current = foundChild;
            if (next == std::wstring::npos) break;
            pos = next + 1;
        }
        return current;
    }
    
    void AddDriveRoots(HWND hTree) {
        DWORD mask = GetLogicalDrives();
        for (int i = 0; i < 26; ++i) {
            if (mask & (1u << i)) {
                wchar_t drive[4] = { wchar_t('A' + i), L':', L'\\', 0 };
                UINT driveType = GetDriveTypeW(drive);
                
                if (driveType == DRIVE_REMOTE) {
                    continue;
                }
                
                std::wstring label = std::wstring(drive);
                bool hasChildren = HasSubdirs(label);
                InsertTreeItem(hTree, TVI_ROOT, label, label, hasChildren);
            }
        }
    }
}
\end{lstlisting}

\textbf{Network}
\begin{lstlisting}[style=CodeListing]
// Winsock must be included before windows.h
#include <winsock2.h>
#include <ws2tcpip.h>
#include <windows.h>
#include <commctrl.h>
#include <winnetwk.h>
#include <iphlpapi.h>
#include "../../include/network.h"
#include "../../include/filesystem.h"
#include <algorithm>

namespace Network {
    bool IsValidDeviceIP(const std::wstring& ip) {
        unsigned int b1, b2, b3, b4;
        if (swscanf(ip.c_str(), L"%u.%u.%u.%u", &b1, &b2, &b3, &b4) != 4) {
            return false;
        }
        
        if (b1 == 0 && b2 == 0 && b3 == 0 && b4 == 0) return false;
        if (b1 == 127) return false;
        if (b1 >= 224 && b1 <= 239) return false;
        if (b1 >= 240) return false;
        if (b1 == 255 && b2 == 255 && b3 == 255 && b4 == 255) return false;
        
        return true;
    }
    
    std::vector<std::wstring> GetNetworkIPsFromARP() {
        std::set<std::wstring> ipSet;
        std::set<std::wstring> localIPs;
        std::wstring defaultGateway;
        
        PMIB_IPADDRTABLE pIpAddrTable = nullptr;
        ULONG ulSize = 0;
        if (GetIpAddrTable(nullptr, &ulSize, FALSE) == ERROR_INSUFFICIENT_BUFFER) {
            pIpAddrTable = (PMIB_IPADDRTABLE)malloc(ulSize);
            if (GetIpAddrTable(pIpAddrTable, &ulSize, FALSE) == NO_ERROR) {
                for (DWORD i = 0; i < pIpAddrTable->dwNumEntries; i++) {
                    MIB_IPADDRROW* row = &pIpAddrTable->table[i];
                    if (row->dwAddr != 0 && row->dwAddr != 0x0100007F) {
                        IN_ADDR addr;
                        addr.S_un.S_addr = row->dwAddr;
                        wchar_t ipStr[16];
                        swprintf(ipStr, 16, L"%d.%d.%d.%d",
                                addr.S_un.S_un_b.s_b1,
                                addr.S_un.S_un_b.s_b2,
                                addr.S_un.S_un_b.s_b3,
                                addr.S_un.S_un_b.s_b4);
                        localIPs.insert(ipStr);
                    }
                }
            }
            free(pIpAddrTable);
        }
        
        PMIB_IPFORWARDTABLE pForwardTable = nullptr;
        ulSize = 0;
        if (GetIpForwardTable(nullptr, &ulSize, FALSE) == ERROR_INSUFFICIENT_BUFFER) {
            pForwardTable = (PMIB_IPFORWARDTABLE)malloc(ulSize);
            if (GetIpForwardTable(pForwardTable, &ulSize, FALSE) == NO_ERROR) {
                for (DWORD i = 0; i < pForwardTable->dwNumEntries; i++) {
                    MIB_IPFORWARDROW* row = &pForwardTable->table[i];
                    if (row->dwForwardDest == 0 && row->dwForwardMask == 0) {
                        IN_ADDR addr;
                        addr.S_un.S_addr = row->dwForwardNextHop;
                        if (addr.S_un.S_addr != 0) {
                            wchar_t gatewayStr[16];
                            swprintf(gatewayStr, 16, L"%d.%d.%d.%d",
                                    addr.S_un.S_un_b.s_b1,
                                    addr.S_un.S_un_b.s_b2,
                                    addr.S_un.S_un_b.s_b3,
                                    addr.S_un.S_un_b.s_b4);
                            defaultGateway = gatewayStr;
                            break;
                        }
                    }
                }
            }
            free(pForwardTable);
        }
        
        PMIB_IPNETTABLE pArpTable = nullptr;
        ulSize = 0;
        if (GetIpNetTable(nullptr, &ulSize, FALSE) == ERROR_INSUFFICIENT_BUFFER) {
            pArpTable = (PMIB_IPNETTABLE)malloc(ulSize);
            if (GetIpNetTable(pArpTable, &ulSize, FALSE) == NO_ERROR) {
                for (DWORD i = 0; i < pArpTable->dwNumEntries; i++) {
                    MIB_IPNETROW* row = &pArpTable->table[i];
                    if (row->dwType == MIB_IPNET_TYPE_DYNAMIC || row->dwType == MIB_IPNET_TYPE_STATIC) {
                        IN_ADDR addr;
                        addr.S_un.S_addr = row->dwAddr;
                        wchar_t ipStr[16];
                        swprintf(ipStr, 16, L"%d.%d.%d.%d",
                                addr.S_un.S_un_b.s_b1,
                                addr.S_un.S_un_b.s_b2,
                                addr.S_un.S_un_b.s_b3,
                                addr.S_un.S_un_b.s_b4);
                        std::wstring ip = ipStr;
                        
                        if (IsValidDeviceIP(ip) && localIPs.find(ip) == localIPs.end()) {
                            ipSet.insert(ip);
                        }
                    }
                }
            }
            free(pArpTable);
        }
        
        if (!defaultGateway.empty() && 
            IsValidDeviceIP(defaultGateway) && 
            localIPs.find(defaultGateway) == localIPs.end() &&
            ipSet.find(defaultGateway) == ipSet.end()) {
            ipSet.insert(defaultGateway);
        }
        
        std::vector<std::wstring> devices;
        for (const auto& ip : ipSet) {
            devices.push_back(L"\\\\" + ip);
        }
        
        return devices;
    }
    
    void EnumerateNetworkResourcesRecursive(NETRESOURCEW* pNetResource, 
                                           std::set<std::wstring>& devicesSet, 
                                           int depth) {
        if (depth > 5) return; // Prevent infinite recursion
        
        HANDLE hEnum = nullptr;
        DWORD dwFlags = 0;
        // Use appropriate flags for enumeration
        if (pNetResource == nullptr) {
            // Root enumeration - find both containers and connectable resources
            dwFlags = RESOURCEUSAGE_CONTAINER | RESOURCEUSAGE_CONNECTABLE;
        } else {
            // Use flags from the resource, but ensure we get containers
            dwFlags = 0; // Will use resource's flags
        }
        
        DWORD result = WNetOpenEnumW(RESOURCE_GLOBALNET, RESOURCETYPE_ANY, dwFlags, pNetResource, &hEnum);
        
        if (result == NO_ERROR && hEnum) {
            DWORD bufferSize = 32768; // Larger buffer for more devices
            NETRESOURCEW* pBuffer = (NETRESOURCEW*)malloc(bufferSize);
            
            while (true) {
                DWORD count = 0xFFFFFFFF;
                result = WNetEnumResourceW(hEnum, &count, pBuffer, &bufferSize);
                
                if (result == NO_ERROR && count > 0) {
                    for (DWORD i = 0; i < count; ++i) {
                        if (pBuffer[i].lpRemoteName) {
                            std::wstring remoteName = pBuffer[i].lpRemoteName;
                            
                            // Add servers (\\computername) to the set
                            if (remoteName.length() >= 2 && remoteName[0] == L'\\' && remoteName[1] == L'\\') {
                                size_t thirdSlash = remoteName.find(L'\\', 2);
                                if (thirdSlash == std::wstring::npos) {
                                    // This is a server name (\\computername), add it
                                    devicesSet.insert(remoteName);
                                } else {
                                    // This is a share (\\computername\share), extract server name
                                    std::wstring serverName = remoteName.substr(0, thirdSlash);
                                    devicesSet.insert(serverName);
                                }
                            }
                            
                            // Recursively enumerate containers (workgroups, domains)
                            if (pBuffer[i].dwUsage & RESOURCEUSAGE_CONTAINER) {
                                EnumerateNetworkResourcesRecursive(&pBuffer[i], devicesSet, depth + 1);
                            }
                        }
                        // Also check for connectable resources that might be servers
                        if ((pBuffer[i].dwUsage & RESOURCEUSAGE_CONNECTABLE) && pBuffer[i].lpRemoteName) {
                            std::wstring remoteName = pBuffer[i].lpRemoteName;
                            if (remoteName.length() >= 2 && remoteName[0] == L'\\' && remoteName[1] == L'\\') {
                                size_t thirdSlash = remoteName.find(L'\\', 2);
                                if (thirdSlash == std::wstring::npos) {
                                    devicesSet.insert(remoteName);
                                } else {
                                    std::wstring serverName = remoteName.substr(0, thirdSlash);
                                    devicesSet.insert(serverName);
                                }
                            }
                        }
                    }
                } else if (result == ERROR_NO_MORE_ITEMS) {
                    break;
                } else {
                    // Error occurred - may need larger buffer or different approach
                    // Try with larger buffer
                    if (result == ERROR_MORE_DATA) {
                        free(pBuffer);
                        bufferSize *= 2;
                        if (bufferSize > 65536) break; // Limit buffer size
                        pBuffer = (NETRESOURCEW*)malloc(bufferSize);
                        continue;
                    }
                    break;
                }
            }
            
            free(pBuffer);
            WNetCloseEnum(hEnum);
        }
    }
    
    std::vector<std::wstring> EnumerateNetworkDrives() {
        std::vector<std::wstring> devices;
        DWORD mask = GetLogicalDrives();
        
        for (int i = 0; i < 26; ++i) {
            if (mask & (1u << i)) {
                wchar_t drive[4] = { wchar_t('A' + i), L':', L'\\', 0 };
                UINT driveType = GetDriveTypeW(drive);
                
                if (driveType == DRIVE_REMOTE) {
                    wchar_t remotePath[MAX_PATH + 1] = {0};
                    DWORD bufSize = MAX_PATH;
                    
                    if (WNetGetConnectionW(drive, remotePath, &bufSize) == NO_ERROR) {
                        devices.push_back(remotePath);
                    } else {
                        devices.push_back(drive);
                    }
                }
            }
        }
        
        return devices;
    }
    
    std::vector<std::wstring> EnumerateUSBDevices() {
        std::vector<std::wstring> devices;
        DWORD mask = GetLogicalDrives();
        
        for (int i = 0; i < 26; ++i) {
            if (mask & (1u << i)) {
                wchar_t drive[4] = { wchar_t('A' + i), L':', L'\\', 0 };
                UINT driveType = GetDriveTypeW(drive);
                
                if (driveType == DRIVE_REMOVABLE || driveType == DRIVE_CDROM || driveType == DRIVE_RAMDISK) {
                    wchar_t volumeName[MAX_PATH + 1] = {0};
                    wchar_t fileSystemName[MAX_PATH + 1] = {0};
                    DWORD serialNumber = 0;
                    
                    if (GetVolumeInformationW(drive, volumeName, MAX_PATH, &serialNumber, nullptr, nullptr, fileSystemName, MAX_PATH)) {
                        devices.push_back(drive);
                    } else {
                        devices.push_back(drive);
                    }
                }
            }
        }
        
        return devices;
    }
    
    std::vector<std::wstring> EnumerateNetworkDevices() {
        std::set<std::wstring> devicesSet;
        std::vector<std::wstring> devices;
        
        // Method 1: Add network drives FIRST (mapped to drive letters like Z:, Y:, etc.)
        // These should always appear in the network view
        std::vector<std::wstring> networkDrives = EnumerateNetworkDrives();
        for (const auto& dev : networkDrives) {
            if (!dev.empty()) {
                devicesSet.insert(dev);
            }
        }
        
        // Method 2: Enumerate all network resources starting from root (like Windows Explorer Network view)
        // This is the main method that finds network computers and shares
        NETRESOURCEW* pNetResource = nullptr;
        EnumerateNetworkResourcesRecursive(pNetResource, devicesSet, 0);
        
        // Method 3: Enumerate workgroups/domains and their servers
        NETRESOURCEW nr{};
        nr.dwScope = RESOURCE_GLOBALNET;
        nr.dwType = RESOURCETYPE_ANY;
        nr.dwUsage = RESOURCEUSAGE_CONTAINER;
        nr.dwDisplayType = RESOURCEDISPLAYTYPE_DOMAIN;
        EnumerateNetworkResourcesRecursive(&nr, devicesSet, 0);
        
        // Method 4: Enumerate servers directly (RESOURCEDISPLAYTYPE_SERVER)
        nr.dwScope = RESOURCE_GLOBALNET;
        nr.dwType = RESOURCETYPE_ANY;
        nr.dwUsage = RESOURCEUSAGE_ALL;
        nr.dwDisplayType = RESOURCEDISPLAYTYPE_SERVER;
        EnumerateNetworkResourcesRecursive(&nr, devicesSet, 0);
        
        // Method 5: Enumerate all disk shares (RESOURCETYPE_DISK)
        nr.dwScope = RESOURCE_GLOBALNET;
        nr.dwType = RESOURCETYPE_DISK;
        nr.dwUsage = RESOURCEUSAGE_ALL;
        nr.dwDisplayType = RESOURCEDISPLAYTYPE_GENERIC;
        EnumerateNetworkResourcesRecursive(&nr, devicesSet, 0);
        
        // Method 6: Try to enumerate from Network Neighborhood root
        // This mimics what Windows Explorer shows in Network view
        HANDLE hEnum = nullptr;
        DWORD result = WNetOpenEnumW(RESOURCE_GLOBALNET, RESOURCETYPE_ANY, 
                                     RESOURCEUSAGE_CONTAINER | RESOURCEUSAGE_CONNECTABLE,
                                     nullptr, &hEnum);
        if (result == NO_ERROR && hEnum) {
            DWORD bufferSize = 16384;
            NETRESOURCEW* pBuffer = (NETRESOURCEW*)malloc(bufferSize);
            DWORD count = 0xFFFFFFFF;
            
            while (true) {
                count = 0xFFFFFFFF;
                result = WNetEnumResourceW(hEnum, &count, pBuffer, &bufferSize);
                if (result == NO_ERROR && count > 0) {
                    for (DWORD i = 0; i < count; i++) {
                        if (pBuffer[i].lpRemoteName) {
                            std::wstring remoteName = pBuffer[i].lpRemoteName;
                            // Add servers (\\computername)
                            if (remoteName.length() >= 2 && remoteName[0] == L'\\' && remoteName[1] == L'\\') {
                                size_t thirdSlash = remoteName.find(L'\\', 2);
                                if (thirdSlash == std::wstring::npos) {
                                    // It's a server, add it
                                    devicesSet.insert(remoteName);
                                }
                            }
                            // Recursively enumerate containers
                            if (pBuffer[i].dwUsage & RESOURCEUSAGE_CONTAINER) {
                                EnumerateNetworkResourcesRecursive(&pBuffer[i], devicesSet, 0);
                            }
                        }
                    }
                } else if (result == ERROR_NO_MORE_ITEMS) {
                    break;
                } else {
                    break;
                }
            }
            free(pBuffer);
            WNetCloseEnum(hEnum);
        }
        
        // Method 7: Add devices from ARP table (IP addresses of devices in local network)
        // This helps find devices that might not be in network enumeration
        std::vector<std::wstring> arpDevices = GetNetworkIPsFromARP();
        for (const auto& dev : arpDevices) {
            if (!dev.empty()) {
                devicesSet.insert(dev);
            }
        }
        
        // Convert set to vector
        for (const auto& dev : devicesSet) {
            if (!dev.empty()) {
                devices.push_back(dev);
            }
        }
        
        return devices;
    }
    
    std::vector<std::wstring> EnumerateNetworkShares(const std::wstring& serverPath) {
        std::vector<std::wstring> shares;
        
        NETRESOURCEW nr{};
        nr.dwScope = RESOURCE_GLOBALNET;
        nr.dwType = RESOURCETYPE_DISK;
        nr.dwUsage = RESOURCEUSAGE_CONNECTABLE;
        nr.dwDisplayType = RESOURCEDISPLAYTYPE_SHARE;
        nr.lpRemoteName = const_cast<wchar_t*>(serverPath.c_str());
        
        HANDLE hEnum = nullptr;
        DWORD result = WNetOpenEnumW(RESOURCE_GLOBALNET, RESOURCETYPE_DISK, 0, &nr, &hEnum);
        
        if (result == NO_ERROR && hEnum) {
            DWORD bufferSize = 16384;
            NETRESOURCEW* pBuffer = (NETRESOURCEW*)malloc(bufferSize);
            DWORD entries = 0xFFFFFFFF;
            
            while (true) {
                DWORD count = entries;
                result = WNetEnumResourceW(hEnum, &count, pBuffer, &bufferSize);
                
                if (result == NO_ERROR) {
                    for (DWORD i = 0; i < count; ++i) {
                        if (pBuffer[i].lpRemoteName) {
                            shares.push_back(pBuffer[i].lpRemoteName);
                        }
                    }
                } else if (result == ERROR_NO_MORE_ITEMS) {
                    break;
                } else {
                    break;
                }
            }
            
            free(pBuffer);
            WNetCloseEnum(hEnum);
        }
        
        return shares;
    }
    
    bool DeviceHasChildren(const std::wstring& device) {
        if (device.length() >= 2 && device[0] == L'\\' && device[1] == L'\\') {
            size_t thirdSlash = device.find(L'\\', 2);
            if (thirdSlash == std::wstring::npos) {
                std::vector<std::wstring> shares = EnumerateNetworkShares(device);
                return !shares.empty();
            } else {
                return Filesystem::HasSubdirs(device);
            }
        } else {
            NETRESOURCEW nr{};
            nr.dwScope = RESOURCE_GLOBALNET;
            nr.dwType = RESOURCETYPE_ANY;
            nr.dwUsage = RESOURCEUSAGE_CONTAINER;
            nr.lpRemoteName = const_cast<wchar_t*>(device.c_str());
            
            HANDLE hEnum = nullptr;
            DWORD enumResult = WNetOpenEnumW(RESOURCE_GLOBALNET, RESOURCETYPE_ANY, 0, &nr, &hEnum);
            if (enumResult == NO_ERROR && hEnum) {
                DWORD bufferSize = 16384;
                NETRESOURCEW* pBuffer = (NETRESOURCEW*)malloc(bufferSize);
                DWORD count = 1;
                enumResult = WNetEnumResourceW(hEnum, &count, pBuffer, &bufferSize);
                bool hasChildren = (enumResult == NO_ERROR && count > 0);
                free(pBuffer);
                WNetCloseEnum(hEnum);
                return hasChildren;
            }
        }
        return false;
    }
    
    void ExpandNetworkDevice(HWND hTree, HTREEITEM hItem) {
        HTREEITEM hChild = TreeView_GetChild(hTree, hItem);
        while (hChild) {
            TVITEMW tvi{}; tvi.mask = TVIF_TEXT; tvi.hItem = hChild;
            wchar_t buf[8]; tvi.pszText = buf; tvi.cchTextMax = 8;
            TreeView_GetItem(hTree, &tvi);
            if (wcscmp(buf, L"...") == 0) {
                TreeView_DeleteItem(hTree, hChild);
                break;
            }
            hChild = TreeView_GetNextSibling(hTree, hChild);
        }
        
        TVITEMW tvi{}; tvi.mask = TVIF_PARAM; tvi.hItem = hItem; TreeView_GetItem(hTree, &tvi);
        Filesystem::NodeData *data = reinterpret_cast<Filesystem::NodeData*>(tvi.lParam);
        if (!data || data->expanded) return;
        
        if (data->path.length() >= 2 && data->path[0] == L'\\' && data->path[1] == L'\\') {
            size_t thirdSlash = data->path.find(L'\\', 2);
            if (thirdSlash == std::wstring::npos) {
                std::vector<std::wstring> shares = EnumerateNetworkShares(data->path);
                
                for (const auto& share : shares) {
                    std::wstring displayName = share;
                    size_t lastSlash = displayName.find_last_of(L'\\');
                    if (lastSlash != std::wstring::npos && lastSlash < displayName.length() - 1) {
                        displayName = displayName.substr(lastSlash + 1);
                    }
                    
                    bool hasChildren = Filesystem::HasSubdirs(share);
                    Filesystem::InsertTreeItem(hTree, hItem, displayName, share, hasChildren);
                }
                data->expanded = true;
                return;
            }
        }
        
        Filesystem::ExpandTreeItem(hTree, hItem);
    }
    
    void AddNetworkDevices(HWND hTree, const std::vector<std::wstring>& devices) {
        // Don't add anything if devices list is empty - let caller handle that
        if (devices.empty()) {
            return;
        }
        
        for (const auto& device : devices) {
            // Skip empty devices
            if (device.empty()) continue;
            
            std::wstring displayName = device;
            
            // Handle UNC paths (\\server\share)
            if (device.length() >= 2 && device[0] == L'\\' && device[1] == L'\\') {
                displayName = device.substr(2);
                // Remove trailing backslashes
                while (!displayName.empty() && displayName.back() == L'\\') {
                    displayName.pop_back();
                }
                // If still empty after removing backslashes, use the original
                if (displayName.empty()) {
                    displayName = device;
                }
            } 
            // Handle drive letters (Z:, Y:, etc.) - network drives
            else if (device.length() >= 3 && device[1] == L':') {
                wchar_t volumeName[MAX_PATH + 1] = {0};
                if (GetVolumeInformationW(device.c_str(), volumeName, MAX_PATH, nullptr, nullptr, nullptr, nullptr, 0)) {
                    if (wcslen(volumeName) > 0) {
                        displayName = std::wstring(device.substr(0, 2)) + L" (" + volumeName + L")";
                    } else {
                        displayName = device.substr(0, 2);
                    }
                } else {
                    displayName = device.substr(0, 2);
                }
            }
            
            if (displayName.empty()) continue;
            
            // Determine if device has children (shares or subdirectories)
            bool hasChildren = false;
            if (device.length() >= 2 && device[0] == L'\\' && device[1] == L'\\') {
                size_t thirdSlash = device.find(L'\\', 2);
                if (thirdSlash == std::wstring::npos) {
                    // This is a server (\\server), check if it has shares
                    hasChildren = DeviceHasChildren(device);
                    // Always assume servers might have children, even if check fails
                    if (!hasChildren) {
                        hasChildren = true; // Allow expansion attempt
                    }
                } else {
                    // This is a share or subdirectory (\\server\share)
                    hasChildren = DeviceHasChildren(device);
                }
            } else if (device.length() >= 3 && device[1] == L':') {
                // Network drive - check if it has subdirectories
                hasChildren = Filesystem::HasSubdirs(device);
            }
            
            Filesystem::InsertTreeItem(hTree, TVI_ROOT, displayName, device, hasChildren);
        }
    }
}
\end{lstlisting}

\textbf{Settings}
\begin{lstlisting}[style=CodeListing]
#include "../../include/settings.h"
#include "../../include/config.h"
#include <algorithm>

namespace Settings {
    // Global state
    static ThemeChoice g_themeChoice = ThemeChoice::System;
    static COLORREF g_colBg = RGB(255,255,255);
    static COLORREF g_colText = RGB(0,0,0);
    static HBRUSH g_hBgBrush = nullptr;
    static std::wstring g_fontFace = L"Segoe UI";
    static int g_fontPt = 12;
    static HFONT g_hFont = nullptr;
    
    ThemeChoice GetTheme() {
        return g_themeChoice;
    }
    
    void SetTheme(ThemeChoice theme) {
        g_themeChoice = theme;
    }
    
    COLORREF GetBackgroundColor() {
        return g_colBg;
    }
    
    COLORREF GetTextColor() {
        return g_colText;
    }
    
    HBRUSH GetBackgroundBrush() {
        return g_hBgBrush;
    }
    
    void SetBackgroundBrush(HBRUSH brush) {
        if (g_hBgBrush && g_hBgBrush != brush) {
            DeleteObject(g_hBgBrush);
        }
        g_hBgBrush = brush;
    }
    
    bool QuerySystemAppsUseLightTheme() {
        DWORD value = 1;
        DWORD size = sizeof(value);
        LSTATUS st = RegGetValueW(HKEY_CURRENT_USER,
            L"Software\\Microsoft\\Windows\\CurrentVersion\\Themes\\Personalize",
            L"AppsUseLightTheme", RRF_RT_REG_DWORD, nullptr, &value, &size);
        if (st != ERROR_SUCCESS) return true; // default light if missing
        return value != 0;
    }
    
    void UpdateThemeColors() {
        ThemeChoice effective = g_themeChoice;
        if (effective == ThemeChoice::System) {
            effective = QuerySystemAppsUseLightTheme() ? ThemeChoice::Light : ThemeChoice::Dark;
        }
        if (effective == ThemeChoice::Dark) {
            g_colBg = RGB(25, 25, 25);
            g_colText = RGB(255, 255, 255);
        } else {
            g_colBg = RGB(255,255,255);
            g_colText = RGB(0,0,0);
        }
        if (g_hBgBrush) { 
            DeleteObject(g_hBgBrush); 
            g_hBgBrush = nullptr; 
        }
        g_hBgBrush = CreateSolidBrush(g_colBg);
    }
    
    std::wstring GetFontFace() {
        return g_fontFace;
    }
    
    void SetFontFace(const std::wstring& face) {
        g_fontFace = face;
    }
    
    int GetFontSize() {
        return g_fontPt;
    }
    
    void SetFontSize(int size) {
        g_fontPt = size;
    }
    
    HFONT GetFont() {
        if (!g_hFont) {
            // Create font if it doesn't exist
            HDC hdc = GetDC(nullptr);
            int logPix = GetDeviceCaps(hdc, LOGPIXELSY);
            ReleaseDC(nullptr, hdc);
            int height = -MulDiv(g_fontPt, logPix, 72);
            g_hFont = CreateFontW(height, 0, 0, 0, FW_NORMAL, FALSE, FALSE, FALSE, DEFAULT_CHARSET,
                                  OUT_DEFAULT_PRECIS, CLIP_DEFAULT_PRECIS, CLEARTYPE_QUALITY,
                                  DEFAULT_PITCH | FF_DONTCARE, g_fontFace.c_str());
        }
        return g_hFont;
    }
    
    void SetFont(HFONT font) {
        if (g_hFont && g_hFont != font) {
            DeleteObject(g_hFont);
        }
        g_hFont = font;
    }
    
    void RecreateFont() {
        if (g_hFont) {
            DeleteObject(g_hFont);
            g_hFont = nullptr;
        }
        GetFont(); // This will create the font
    }
    
    void Save() {
        HKEY hKey;
        if (RegCreateKeyExW(HKEY_CURRENT_USER, L"Software\\KursWork\\FileManager", 0, nullptr, 0, KEY_SET_VALUE, nullptr, &hKey, nullptr) == ERROR_SUCCESS) {
            DWORD theme = (DWORD)g_themeChoice;
            RegSetValueExW(hKey, L"Theme", 0, REG_DWORD, reinterpret_cast<const BYTE*>(&theme), sizeof(theme));
            DWORD size = (DWORD)g_fontPt;
            RegSetValueExW(hKey, L"FontSize", 0, REG_DWORD, reinterpret_cast<const BYTE*>(&size), sizeof(size));
            RegSetValueExW(hKey, L"FontFace", 0, REG_SZ, reinterpret_cast<const BYTE*>(g_fontFace.c_str()), (DWORD)((g_fontFace.size() + 1) * sizeof(wchar_t)));
            RegCloseKey(hKey);
        }
    }
    
    void Load() {
        HKEY hKey;
        if (RegOpenKeyExW(HKEY_CURRENT_USER, L"Software\\KursWork\\FileManager", 0, KEY_QUERY_VALUE, &hKey) == ERROR_SUCCESS) {
            DWORD theme = (DWORD)ThemeChoice::System; 
            DWORD sz = sizeof(theme);
            if (RegGetValueW(hKey, nullptr, L"Theme", RRF_RT_REG_DWORD, nullptr, &theme, &sz) == ERROR_SUCCESS) {
                if (theme <= (DWORD)ThemeChoice::Dark) g_themeChoice = (ThemeChoice)theme;
            }
            DWORD size = 12; 
            sz = sizeof(size);
            if (RegGetValueW(hKey, nullptr, L"FontSize", RRF_RT_REG_DWORD, nullptr, &size, &sz) == ERROR_SUCCESS) {
                if (size >= 8 && size <= 72) g_fontPt = (int)size;
            }
            wchar_t buf[128]; 
            DWORD bytes = sizeof(buf);
            if (RegGetValueW(hKey, nullptr, L"FontFace", RRF_RT_REG_SZ, nullptr, buf, &bytes) == ERROR_SUCCESS) {
                g_fontFace = buf;
            }
            RegCloseKey(hKey);
        }
        UpdateThemeColors();
    }
}
\end{lstlisting}

\textbf{UI}
\begin{lstlisting}[style=CodeListing]
// Control initialization module

#ifndef UNICODE
#define UNICODE
#endif
#ifndef _UNICODE
#define _UNICODE
#endif

#include "../../include/ui.h"
#include "../../include/filesystem.h"
#include "../../include/network.h"
#include "../../include/settings.h"
#include "../../include/config.h"
#include <windows.h>
#include <commctrl.h>
#include <shellapi.h>
#include <gdiplus.h>
#include <thread>
#include <mutex>

using namespace Gdiplus;

// Forward declarations
extern WindowData* GetWindowData(HWND hwnd);
extern void SetWindowDataInternal(HWND hwnd, WindowData* data);
extern LRESULT CALLBACK BorderSubclassProc(HWND hwnd, UINT uMsg, WPARAM wParam, LPARAM lParam, UINT_PTR uIdSubclass, DWORD_PTR dwRefData);
extern LRESULT CALLBACK ListViewSubclassProc(HWND hwnd, UINT uMsg, WPARAM wParam, LPARAM lParam);

namespace UI {
    void AddTooltipFor(HWND hParent, HWND hCtrl, const wchar_t *text) {
        HWND hToolTip = GetToolTip(hParent);
        if (!hToolTip) {
            hToolTip = CreateWindowExW(WS_EX_TOPMOST, TOOLTIPS_CLASSW, nullptr,
                WS_POPUP | TTS_NOPREFIX | TTS_ALWAYSTIP,
                CW_USEDEFAULT, CW_USEDEFAULT, CW_USEDEFAULT, CW_USEDEFAULT,
                hParent, nullptr, GetModuleHandleW(nullptr), nullptr);
            SetWindowPos(hToolTip, HWND_TOPMOST, 0, 0, 0, 0,
                SWP_NOMOVE | SWP_NOSIZE | SWP_NOACTIVATE);
            SendMessageW(hToolTip, TTM_SETMAXTIPWIDTH, 0, 300);
            SetToolTip(hParent, hToolTip);
        }
        TOOLINFOW ti{}; ti.cbSize = sizeof(ti);
        ti.uFlags = TTF_SUBCLASS | TTF_IDISHWND;
        ti.hwnd = hParent;
        ti.uId = (UINT_PTR)hCtrl;
        ti.lpszText = const_cast<wchar_t*>(text);
        GetClientRect(hCtrl, &ti.rect);
        SendMessageW(hToolTip, TTM_ADDTOOL, 0, (LPARAM)&ti);
    }
    
    HICON LoadIconFromPNG(const wchar_t *path, ULONG_PTR gdiplusToken) {
        HICON hIcon = nullptr;
        Image *image = new Image(path);
        if (image && image->GetLastStatus() == Ok) {
            int smallIconSize = GetSystemMetrics(SM_CXSMICON);
            int largeIconSize = GetSystemMetrics(SM_CXICON);
            int iconSize = std::max(largeIconSize, 32);
            if (iconSize < 16) iconSize = 32;
            if (iconSize > 256) iconSize = 256;
            Bitmap *bitmap = new Bitmap(iconSize, iconSize);
            Graphics *graphics = Graphics::FromImage(bitmap);
            graphics->SetInterpolationMode(InterpolationModeHighQualityBicubic);
            graphics->SetSmoothingMode(SmoothingModeHighQuality);
            graphics->SetPixelOffsetMode(PixelOffsetModeHighQuality);
            graphics->DrawImage(image, 0, 0, iconSize, iconSize);
            bitmap->GetHICON(&hIcon);
            delete graphics;
            delete bitmap;
            delete image;
        } else {
            if (image) delete image;
        }
        return hIcon;
    }
    
    void InitializeControls(HWND hwnd, UIControls &controls, UIState &state) {
        INITCOMMONCONTROLSEX icc{ sizeof(icc), ICC_TREEVIEW_CLASSES | ICC_LISTVIEW_CLASSES };
        InitCommonControlsEx(&icc);
        
        UI::BuildMenu(hwnd);
        UI::BuildAccelerators(hwnd, controls.hAccel);
        Settings::Load();
        
        // Create and store WindowData
        WindowData* wndData = UI::CreateWindowData();
        UI::InitializeWindowData(wndData, controls, state);
        SetWindowDataInternal(hwnd, wndData);
        
        // Create buttons with Unicode arrow symbols
        controls.hBackBtn = CreateWindowExW(0, L"BUTTON", L"◀", 
                                             WS_CHILD | WS_VISIBLE | WS_TABSTOP | BS_PUSHBUTTON,
                                             0, 0, 28, 28, hwnd, (HMENU)ID_BTN_BACK, GetModuleHandleW(nullptr), nullptr);
        controls.hFwdBtn = CreateWindowExW(0, L"BUTTON", L"◀",
                                             WS_CHILD | WS_VISIBLE | WS_TABSTOP | BS_PUSHBUTTON,
                                             0, 0, 28, 28, hwnd, (HMENU)ID_BTN_FWD, GetModuleHandleW(nullptr), nullptr);
        SetWindowTextW(controls.hFwdBtn, L"▶");
        controls.hViewDetailsBtn = CreateWindowExW(0, L"BUTTON", L"☰", 
                                             WS_CHILD | WS_VISIBLE | WS_TABSTOP | BS_PUSHBUTTON,
                                             0, 0, 28, 28, hwnd, (HMENU)ID_BTN_VIEW_DETAILS, GetModuleHandleW(nullptr), nullptr);
        controls.hViewLargeBtn = CreateWindowExW(0, L"BUTTON", L"⬛",
                                             WS_CHILD | WS_VISIBLE | WS_TABSTOP | BS_PUSHBUTTON,
                                             0, 0, 28, 28, hwnd, (HMENU)ID_BTN_VIEW_LARGE, GetModuleHandleW(nullptr), nullptr);
        controls.hViewSmallBtn = CreateWindowExW(0, L"BUTTON", L"▤",
                                             WS_CHILD | WS_VISIBLE | WS_TABSTOP | BS_PUSHBUTTON,
                                             0, 0, 28, 28, hwnd, (HMENU)ID_BTN_VIEW_SMALL, GetModuleHandleW(nullptr), nullptr);
        
        // Create path edit field
        controls.hPathEdit = CreateWindowExW(0, L"EDIT", L"",
                                          WS_CHILD | WS_VISIBLE | WS_TABSTOP | ES_AUTOHSCROLL | WS_BORDER,
                                          0, 0, 200, 24, hwnd, (HMENU)ID_EDIT_PATH, GetModuleHandleW(nullptr), nullptr);
        SetWindowSubclass(controls.hPathEdit, BorderSubclassProc, 1, 0);
        
        // Create Go button
        controls.hGoBtn = CreateWindowExW(0, L"BUTTON", L"Перейти", WS_CHILD | WS_VISIBLE | WS_TABSTOP,
                                       0, 0, 60, 24, hwnd, (HMENU)ID_BTN_GO, GetModuleHandleW(nullptr), nullptr);
        
        // Create trees
        controls.hTree = CreateWindowExW(0, WC_TREEVIEWW, L"",
                                      WS_CHILD | WS_VISIBLE | TVS_HASLINES | TVS_LINESATROOT | TVS_HASBUTTONS | WS_BORDER,
                                      0, 0, 100, 100, hwnd, (HMENU)ID_TREE_TOP, GetModuleHandleW(nullptr), nullptr);
        SetWindowSubclass(controls.hTree, BorderSubclassProc, 1, 0);
        
        // Create splitter
        controls.hSplitter = CreateWindowExW(0, L"STATIC", L"",
                                        WS_CHILD | WS_VISIBLE,
                                        0, 0, 100, 4, hwnd, (HMENU)ID_SPLITTER, GetModuleHandleW(nullptr), nullptr);
        
        // Create bottom tree for network devices
        controls.hTreeBottom = CreateWindowExW(0, WC_TREEVIEWW, L"",
                                           WS_CHILD | WS_VISIBLE | TVS_HASLINES | TVS_LINESATROOT | TVS_HASBUTTONS | WS_BORDER,
                                           0, 0, 100, 100, hwnd, (HMENU)ID_TREE_BOTTOM, GetModuleHandleW(nullptr), nullptr);
        SetWindowSubclass(controls.hTreeBottom, BorderSubclassProc, 1, 0);
        
        // Create list view
        controls.hList = CreateWindowExW(0, WC_LISTVIEWW, L"",
                                      WS_CHILD | WS_VISIBLE | LVS_REPORT | LVS_SINGLESEL | LVS_SHAREIMAGELISTS | WS_BORDER,
                                      0, 0, 100, 100, hwnd, (HMENU)ID_LIST, GetModuleHandleW(nullptr), nullptr);
        SetWindowSubclass(controls.hList, BorderSubclassProc, 1, 0);
        
        // Attach system imagelists for icons
        {
            SHFILEINFOW sfi{};
            HIMAGELIST hSmall = (HIMAGELIST)SHGetFileInfoW(L"C:\\", FILE_ATTRIBUTE_DIRECTORY, &sfi, sizeof(sfi), SHGFI_SYSICONINDEX | SHGFI_SMALLICON);
            HIMAGELIST hLarge = (HIMAGELIST)SHGetFileInfoW(L"C:\\", FILE_ATTRIBUTE_DIRECTORY, &sfi, sizeof(sfi), SHGFI_SYSICONINDEX | SHGFI_LARGEICON);
            if (hSmall) ListView_SetImageList(controls.hList, hSmall, LVSIL_SMALL);
            if (hLarge) ListView_SetImageList(controls.hList, hLarge, LVSIL_NORMAL);
        }
        
        // Set up list view columns
        SetListColumns(controls.hList);
        AdjustReportColumns(controls.hList);
        DWORD ex = ListView_GetExtendedListViewStyle(controls.hList);
        ex |= LVS_EX_FULLROWSELECT | LVS_EX_LABELTIP;
        ThemeChoice effective = Settings::GetTheme();
        if (effective == ThemeChoice::System) {
            effective = Settings::QuerySystemAppsUseLightTheme() ? ThemeChoice::Light : ThemeChoice::Dark;
        }
        if (effective != ThemeChoice::Dark) {
            ex |= LVS_EX_GRIDLINES;
        }
        ListView_SetExtendedListViewStyle(controls.hList, ex);
        
        // Subclass ListView to draw empty message
        state.oldListProc = (WNDPROC)SetWindowLongPtrW(controls.hList, GWLP_WNDPROC, (LONG_PTR)ListViewSubclassProc);
        
        // Create empty label
        controls.hEmptyLabel = CreateWindowExW(0, L"STATIC", L"Выберите директорию",
                                           WS_CHILD | SS_CENTER | SS_CENTERIMAGE,
                                           0, 0, 100, 100, hwnd, nullptr, GetModuleHandleW(nullptr), nullptr);
        ShowWindow(controls.hEmptyLabel, SW_HIDE);
        
        // Create status bar
        controls.hStatusBar = CreateWindowExW(0, STATUSCLASSNAMEW, L"",
                                             WS_CHILD | WS_VISIBLE | SBARS_SIZEGRIP,
                                             0, 0, 0, 0, hwnd, (HMENU)ID_STATUS_BAR, GetModuleHandleW(nullptr), nullptr);
        int parts[] = { -1 };
        SendMessageW(controls.hStatusBar, SB_SETPARTS, 1, (LPARAM)parts);
        
        AdjustReportColumns(controls.hList);
        UpdateListViewRowHeight(controls.hList);
        Filesystem::AddDriveRoots(controls.hTree);
        
        // Clear bottom tree and add loading placeholder
        TreeView_DeleteAllItems(controls.hTreeBottom);
        state.loadingItem = Filesystem::InsertTreeItem(controls.hTreeBottom, TVI_ROOT, L"Загрузка...", L"", false);
        UI::UpdateWindowState(hwnd, state);
        
        LayoutChildren(hwnd, controls, state);
        
        // Start network scan in background thread
        state.networkScanInProgress = true;
        UI::UpdateWindowState(hwnd, state);
        std::thread([hwnd]() {
            std::vector<std::wstring> devices;
            try {
                Sleep(100);
                devices = Network::EnumerateNetworkDevices();
            } catch (const std::exception& e) {
                devices.clear();
            } catch (...) {
                devices.clear();
            }
            
            UI::UpdateScannedDevices(hwnd, devices);
            PostMessageW(hwnd, WM_USER_NETWORK_SCAN_COMPLETE, 0, 0);
        }).detach();
        
        UI::UpdateWindowData(hwnd, controls, state);
        
        EnableWindow(controls.hBackBtn, FALSE);
        EnableWindow(controls.hFwdBtn, FALSE);
        UI::ApplyThemeToControls(hwnd, controls);
        UI::ApplyFontToControls(hwnd, controls);
        UI::UpdateMenuChecks(hwnd);
        
        if (state.currentPath.empty()) {
            ShowMessage(controls.hList, controls.hEmptyLabel, L"Выберите директорию");
        }
        
        UI::AddTooltipFor(hwnd, controls.hBackBtn, L"Назад\nAlt+Left");
        UI::AddTooltipFor(hwnd, controls.hFwdBtn, L"Вперёд\nAlt+Right");
        UI::AddTooltipFor(hwnd, controls.hViewDetailsBtn, L"Вид: Подробно");
        UI::AddTooltipFor(hwnd, controls.hViewLargeBtn, L"Вид: Крупные значки");
        UI::AddTooltipFor(hwnd, controls.hViewSmallBtn, L"Вид: Мелкие значки");
        if (controls.hGoBtn) UI::AddTooltipFor(hwnd, controls.hGoBtn, L"Перейти к указанному пути");
    }
}

// Layout management module

#ifndef UNICODE
#define UNICODE
#endif
#ifndef _UNICODE
#define _UNICODE
#endif

#include "../../include/ui.h"
#include "../../include/settings.h"
#include <windows.h>

namespace UI {
    void LayoutChildren(HWND hwnd, UIControls &controls, UIState &state) {
        RECT rc; GetClientRect(hwnd, &rc);
        int width = rc.right - rc.left;
        int height = rc.bottom - rc.top;
        int leftW = width / 3;
        int rightX = leftW;
        
        // Calculate sizes based on font size for adaptive layout
        int fontSize = Settings::GetFontSize();
        HDC hdc = GetDC(hwnd);
        int logPix = GetDeviceCaps(hdc, LOGPIXELSY);
        ReleaseDC(hwnd, hdc);
        // Base size scales with font (12pt = 32px control bar, scale proportionally)
        int ctlBarH = MulDiv(fontSize * 8, logPix, 72 * 3); // Approximately fontSize * 2.67
        if (ctlBarH < 24) ctlBarH = 24; // Minimum height
        if (ctlBarH > 60) ctlBarH = 60; // Maximum height
        
        int pad = fontSize / 2; // Padding scales with font
        if (pad < 4) pad = 4;
        if (pad > 12) pad = 12;
        
        int btnH = ctlBarH - pad - pad/2;
        int btnW = btnH; // Square buttons
        int x = rightX + pad;
        int splitterH = 4;
        
        int minTreeH = 50;
        if (state.splitterPos < minTreeH) state.splitterPos = minTreeH;
        if (state.splitterPos > height - minTreeH - splitterH) state.splitterPos = height - minTreeH - splitterH;
        
        int topTreeH = state.splitterPos;
        MoveWindow(controls.hTree, 0, 0, leftW, topTreeH, TRUE);
        MoveWindow(controls.hSplitter, 0, topTreeH, leftW, splitterH, TRUE);
        
        int bottomTreeY = topTreeH + splitterH;
        int bottomTreeH = height - bottomTreeY;
        MoveWindow(controls.hTreeBottom, 0, bottomTreeY, leftW, bottomTreeH, TRUE);
        
        MoveWindow(controls.hBackBtn, x, pad, btnW, btnH, TRUE); x += btnW + pad;
        MoveWindow(controls.hFwdBtn, x, pad, btnW, btnH, TRUE); x += btnW + pad*2;
        MoveWindow(controls.hViewDetailsBtn, x, pad, btnW, btnH, TRUE); x += btnW + pad;
        MoveWindow(controls.hViewLargeBtn, x, pad, btnW, btnH, TRUE); x += btnW + pad;
        MoveWindow(controls.hViewSmallBtn, x, pad, btnW, btnH, TRUE); x += btnW + pad*2;
        
        int pathEditX = x;
        // Go button width scales with font
        int goBtnW = MulDiv(fontSize * 5, logPix, 72); // Approximately fontSize * 0.7
        if (goBtnW < 50) goBtnW = 50;
        if (goBtnW > 100) goBtnW = 100;
        int pathEditW = width - pathEditX - goBtnW - pad;
        if (controls.hPathEdit && controls.hGoBtn) {
            if (pathEditW > 100) {
                ShowWindow(controls.hPathEdit, SW_SHOW);
                MoveWindow(controls.hPathEdit, pathEditX, pad, pathEditW, btnH, TRUE);
                MoveWindow(controls.hGoBtn, pathEditX + pathEditW, pad, goBtnW, btnH, TRUE);
            } else {
                ShowWindow(controls.hPathEdit, SW_HIDE);
                MoveWindow(controls.hGoBtn, pathEditX, pad, goBtnW, btnH, TRUE);
            }
        }
        
        int listX = rightX;
        int listY = ctlBarH;
        int listW = width - rightX;
        int listH = height - ctlBarH;
        MoveWindow(controls.hList, listX, listY, listW, listH, TRUE);
        if (controls.hEmptyLabel) {
            MoveWindow(controls.hEmptyLabel, listX, listY, listW, listH, TRUE);
        }
        if (controls.hStatusBar) {
            RECT statusRect;
            GetClientRect(controls.hStatusBar, &statusRect);
            MoveWindow(controls.hStatusBar, listX, height - statusRect.bottom, listW, statusRect.bottom, TRUE);
        }
    }
}

// List view operations module

#ifndef UNICODE
#define UNICODE
#endif
#ifndef _UNICODE
#define _UNICODE
#endif

#include "../../include/ui.h"
#include "../../include/settings.h"
#include <windows.h>
#include <commctrl.h>
#include <shellapi.h>

namespace UI {
    // Helper function to get icon size based on font size
    int GetIconSize() {
        int fontSize = Settings::GetFontSize();
        HDC hdc = GetDC(nullptr);
        int logPix = GetDeviceCaps(hdc, LOGPIXELSY);
        ReleaseDC(nullptr, hdc);
        // Icon size should be roughly 1.5x font height for good visibility
        int iconSize = MulDiv(fontSize * 3, logPix, 72 * 2); // fontSize * 1.5
        // Clamp to reasonable sizes (16-48 pixels)
        if (iconSize < 16) iconSize = 16;
        if (iconSize > 48) iconSize = 48;
        return iconSize;
    }
    
    // Update list view row height to accommodate proportional icon size
    void UpdateListViewRowHeight(HWND hList) {
        int fontSize = Settings::GetFontSize();
        HDC hdc = GetDC(nullptr);
        int logPix = GetDeviceCaps(hdc, LOGPIXELSY);
        ReleaseDC(nullptr, hdc);
        // Row height should be proportional to font size (1.8x font height for icon + text)
        int rowHeight = MulDiv(fontSize * 9, logPix, 72 * 5); // fontSize * 1.8
        // Clamp to reasonable sizes
        if (rowHeight < 20) rowHeight = 20;
        if (rowHeight > 60) rowHeight = 60;
        // Use LVM_SETITEMHEIGHT message directly
        #ifndef LVM_FIRST
        #define LVM_FIRST 0x1000
        #endif
        #ifndef LVM_SETITEMHEIGHT
        #define LVM_SETITEMHEIGHT (LVM_FIRST + 141)
        #endif
        SendMessageW(hList, LVM_SETITEMHEIGHT, 0, (LPARAM)rowHeight);
    }
    
    void SetListColumns(HWND hList) {
        // Set extended style - gridlines will be drawn manually in dark theme
        ThemeChoice effective = Settings::GetTheme();
        if (effective == ThemeChoice::System) {
            effective = Settings::QuerySystemAppsUseLightTheme() ? ThemeChoice::Light : ThemeChoice::Dark;
        }
        DWORD exStyle = LVS_EX_FULLROWSELECT | LVS_EX_LABELTIP;
        if (effective != ThemeChoice::Dark) {
            exStyle |= LVS_EX_GRIDLINES; // Only use system gridlines in light theme
        }
        ListView_SetExtendedListViewStyle(hList, exStyle);
        
        LVCOLUMNW col{};
        col.mask = LVCF_TEXT | LVCF_WIDTH | LVCF_SUBITEM | LVCF_FMT;
        col.fmt = LVCFMT_LEFT;
        wchar_t nameCol[] = L"Name";
        wchar_t typeCol[] = L"Type";
        wchar_t modCol[] = L"Modified";
        col.pszText = nameCol;
        col.cx = 360; col.iSubItem = 0; 
        ListView_InsertColumn(hList, 0, &col);
        col.pszText = typeCol;
        col.cx = 120; col.iSubItem = 1; 
        ListView_InsertColumn(hList, 1, &col);
        col.pszText = modCol;
        col.cx = 200; col.iSubItem = 2; 
        ListView_InsertColumn(hList, 2, &col);
    }
    
    void AdjustReportColumns(HWND hList) {
        LONG_PTR style = GetWindowLongPtrW(hList, GWL_STYLE);
        if ((style & LVS_TYPEMASK) != LVS_REPORT) return;
        RECT rc{}; GetClientRect(hList, &rc);
        int clientW = rc.right - rc.left;
        if (GetWindowLongPtrW(hList, GWL_STYLE) & WS_VSCROLL) {
            clientW -= GetSystemMetrics(SM_CXVSCROLL);
        }
        if (clientW <= 0) return;
        int typeW = 120;
        int modW = 200;
        int minCol = 60;
        int nameW = clientW - typeW - modW;
        if (nameW < minCol) {
            int deficit = minCol - nameW;
            nameW = minCol;
            int reduceType = std::min(deficit / 2, std::max(0, typeW - minCol));
            typeW -= reduceType;
            deficit -= reduceType;
            int reduceMod = std::min(deficit, std::max(0, modW - minCol));
            modW -= reduceMod;
            deficit -= reduceMod;
            if (deficit > 0) nameW = std::max(minCol, nameW - deficit);
        }
        ListView_SetColumnWidth(hList, 1, typeW);
        ListView_SetColumnWidth(hList, 2, modW);
        ListView_SetColumnWidth(hList, 0, std::max(1, clientW - typeW - modW));
        ShowScrollBar(hList, SB_HORZ, FALSE);
    }
    
    void SetListViewMode(HWND hList, DWORD mode) {
        LONG_PTR style = GetWindowLongPtrW(hList, GWL_STYLE);
        style &= ~(LVS_TYPEMASK);
        style |= mode;
        SetWindowLongPtrW(hList, GWL_STYLE, style);
        if ((mode & LVS_TYPEMASK) == LVS_REPORT) {
            while (ListView_DeleteColumn(hList, 0)) {}
            SetListColumns(hList);
            DWORD ex = ListView_GetExtendedListViewStyle(hList);
            ex |= LVS_EX_FULLROWSELECT | LVS_EX_LABELTIP;
            ListView_SetExtendedListViewStyle(hList, ex);
            // Ensure image lists are set for report mode too (icons in first column)
            SHFILEINFOW sfi{};
            HIMAGELIST hSmall = (HIMAGELIST)SHGetFileInfoW(L"C:\\", FILE_ATTRIBUTE_DIRECTORY, &sfi, sizeof(sfi), SHGFI_SYSICONINDEX | SHGFI_SMALLICON);
            if (hSmall) {
                ListView_SetImageList(hList, hSmall, LVSIL_SMALL);
            }
        }
        // Ensure image lists are set for icon modes
        if ((mode & LVS_TYPEMASK) == LVS_ICON || (mode & LVS_TYPEMASK) == LVS_SMALLICON) {
            SHFILEINFOW sfi{};
            HIMAGELIST hSmall = (HIMAGELIST)SHGetFileInfoW(L"C:\\", FILE_ATTRIBUTE_DIRECTORY, &sfi, sizeof(sfi), SHGFI_SYSICONINDEX | SHGFI_SMALLICON);
            HIMAGELIST hLarge = (HIMAGELIST)SHGetFileInfoW(L"C:\\", FILE_ATTRIBUTE_DIRECTORY, &sfi, sizeof(sfi), SHGFI_SYSICONINDEX | SHGFI_LARGEICON);
            if (hSmall) ListView_SetImageList(hList, hSmall, LVSIL_SMALL);
            if (hLarge) ListView_SetImageList(hList, hLarge, LVSIL_NORMAL);
        }
        InvalidateRect(hList, nullptr, TRUE);
    }
    
    bool IsReportMode(HWND hList) {
        LONG_PTR style = GetWindowLongPtrW(hList, GWL_STYLE);
        return (style & LVS_TYPEMASK) == LVS_REPORT;
    }
    
    bool IsIconMode(HWND hList) {
        LONG_PTR style = GetWindowLongPtrW(hList, GWL_STYLE);
        DWORD t = (DWORD)(style & LVS_TYPEMASK);
        return t == LVS_ICON || t == LVS_SMALLICON;
    }
    
    void AdjustIconLayout(HWND hList) {
        if (!IsIconMode(hList)) return;
        RECT rc{}; GetClientRect(hList, &rc);
        int clientW = rc.right - rc.left;
        if (clientW <= 0) return;
        LONG_PTR style = GetWindowLongPtrW(hList, GWL_STYLE);
        BOOL isSmall = ((style & LVS_TYPEMASK) == LVS_SMALLICON);
        if (isSmall) {
            style &= ~LVS_ALIGNLEFT;
            style |= LVS_AUTOARRANGE | LVS_ALIGNTOP;
        } else {
            style &= ~LVS_ALIGNTOP;
            style |= LVS_AUTOARRANGE | LVS_ALIGNLEFT;
        }
        SetWindowLongPtrW(hList, GWL_STYLE, style);
        LRESULT spacing = SendMessageW(hList, LVM_GETITEMSPACING, (WPARAM)isSmall, 0);
        int defX = LOWORD(spacing);
        int defY = HIWORD(spacing);
        if (defX <= 0) defX = 100;
        int cols = std::max(1, clientW / defX);
        int newX = std::max(32, clientW / cols);
        SendMessageW(hList, LVM_SETICONSPACING, 0, MAKELPARAM(newX, defY));
        ListView_Arrange(hList, LVA_DEFAULT);
        ShowScrollBar(hList, SB_HORZ, FALSE);
    }
    
    void ClearList(HWND hList) {
        ListView_DeleteAllItems(hList);
    }
    
    void ShowMessage(HWND hList, HWND hEmptyLabel, const std::wstring &message) {
        ClearList(hList);
        
        // Hide list view header when showing message
        HWND hHeader = ListView_GetHeader(hList);
        if (hHeader) {
            ShowWindow(hHeader, SW_HIDE);
        }
        
        if (hEmptyLabel) {
            SetWindowTextW(hEmptyLabel, message.c_str());
            ShowWindow(hEmptyLabel, SW_SHOW);
            SetWindowPos(hEmptyLabel, HWND_TOP, 0, 0, 0, 0, SWP_NOMOVE | SWP_NOSIZE);
            InvalidateRect(hEmptyLabel, nullptr, TRUE);
        }
        ShowWindow(hList, SW_HIDE);
        InvalidateRect(hList, nullptr, TRUE);
        UpdateWindow(hList);
        RedrawWindow(hList, nullptr, nullptr, RDW_INVALIDATE | RDW_UPDATENOW | RDW_ERASE);
    }
    
    void UpdateStatusBar(HWND hStatusBar, const std::vector<Filesystem::ListItem> &items) {
        if (!hStatusBar) return;
        
        int fileCount = 0;
        int dirCount = 0;
        __int64 totalSize = 0;
        
        for (const auto& item : items) {
            if (item.isDir) {
                dirCount++;
            } else {
                fileCount++;
                WIN32_FIND_DATAW fd{};
                HANDLE h = FindFirstFileW(item.fullPath.c_str(), &fd);
                if (h != INVALID_HANDLE_VALUE) {
                    totalSize += ((__int64)fd.nFileSizeHigh << 32) | fd.nFileSizeLow;
                    FindClose(h);
                }
            }
        }
        
        wchar_t statusText[256];
        if (fileCount == 0 && dirCount == 0) {
            swprintf(statusText, 256, L"Готово");
        } else {
            swprintf(statusText, 256, L"Файлов: %d, Папок: %d", fileCount, dirCount);
        }
        SendMessageW(hStatusBar, SB_SETTEXT, 0, (LPARAM)statusText);
    }
}

// Menu and accelerators module

#ifndef UNICODE
#define UNICODE
#endif
#ifndef _UNICODE
#define _UNICODE
#endif

#include "../../include/ui.h"
#include "../../include/settings.h"
#include "../../include/config.h"
#include <windows.h>

namespace UI {
    void BuildMenu(HWND hwnd) {
        HMENU hBar = CreateMenu();
        HMENU hFile = CreatePopupMenu();
        HMENU hSettings = CreatePopupMenu();
        HMENU hTheme = CreatePopupMenu();
        HMENU hFont = CreatePopupMenu();
        HMENU hFontSize = CreatePopupMenu();
        
        AppendMenuW(hFile, MF_STRING, ID_MENU_FILE_ABOUT, L"О программе\tF1");
        AppendMenuW(hFile, MF_SEPARATOR, 0, nullptr);
        AppendMenuW(hFile, MF_STRING, ID_MENU_FILE_EXIT, L"Выйти\tCtrl+Q");
        
        AppendMenuW(hTheme, MF_STRING | MF_CHECKED, ID_MENU_THEME_SYSTEM, L"Системная\tCtrl+1");
        AppendMenuW(hTheme, MF_STRING, ID_MENU_THEME_LIGHT, L"Светлая\tCtrl+2");
        AppendMenuW(hTheme, MF_STRING, ID_MENU_THEME_DARK, L"Тёмная\tCtrl+3");
        
        AppendMenuW(hFont, MF_STRING | MF_CHECKED, ID_MENU_FONT_SEGOE, L"Segoe UI\tCtrl+4");
        AppendMenuW(hFont, MF_STRING, ID_MENU_FONT_CONSOLAS, L"Consolas\tCtrl+5");
        AppendMenuW(hFont, MF_STRING, ID_MENU_FONT_COURIER, L"Courier New\tCtrl+6");
        
        AppendMenuW(hFontSize, MF_STRING, ID_MENU_FONTSIZE_10, L"10\tCtrl+Shift+1");
        AppendMenuW(hFontSize, MF_STRING | MF_CHECKED, ID_MENU_FONTSIZE_12, L"12\tCtrl+Shift+2");
        AppendMenuW(hFontSize, MF_STRING, ID_MENU_FONTSIZE_14, L"14\tCtrl+Shift+3");
        AppendMenuW(hFontSize, MF_STRING, ID_MENU_FONTSIZE_16, L"16\tCtrl+Shift+4");
        
        AppendMenuW(hSettings, MF_POPUP, (UINT_PTR)hTheme, L"Тема");
        AppendMenuW(hSettings, MF_POPUP, (UINT_PTR)hFont, L"Шрифт");
        AppendMenuW(hSettings, MF_POPUP, (UINT_PTR)hFontSize, L"Размер текста");
        
        AppendMenuW(hBar, MF_POPUP, (UINT_PTR)hFile, L"Файл");
        AppendMenuW(hBar, MF_POPUP, (UINT_PTR)hSettings, L"Настройки");
        
        SetMenu(hwnd, hBar);
    }
    
    void BuildAccelerators(HWND hwnd, HACCEL &hAccel) {
        ACCEL acc[] = {
            { FVIRTKEY, VK_F1, ID_MENU_FILE_ABOUT },
            { FVIRTKEY | FCONTROL, 'Q', ID_MENU_FILE_EXIT },
            { FVIRTKEY | FCONTROL, '1', ID_MENU_THEME_SYSTEM },
            { FVIRTKEY | FCONTROL, '2', ID_MENU_THEME_LIGHT },
            { FVIRTKEY | FCONTROL, '3', ID_MENU_THEME_DARK },
            { FVIRTKEY | FCONTROL, '4', ID_MENU_FONT_SEGOE },
            { FVIRTKEY | FCONTROL, '5', ID_MENU_FONT_CONSOLAS },
            { FVIRTKEY | FCONTROL, '6', ID_MENU_FONT_COURIER },
            { FVIRTKEY | FCONTROL | FSHIFT, '1', ID_MENU_FONTSIZE_10 },
            { FVIRTKEY | FCONTROL | FSHIFT, '2', ID_MENU_FONTSIZE_12 },
            { FVIRTKEY | FCONTROL | FSHIFT, '3', ID_MENU_FONTSIZE_14 },
            { FVIRTKEY | FCONTROL | FSHIFT, '4', ID_MENU_FONTSIZE_16 },
            { FVIRTKEY | FCONTROL, VK_OEM_MINUS, ID_CMD_FONTSIZE_DEC },
            { FVIRTKEY | FCONTROL, VK_OEM_PLUS, ID_CMD_FONTSIZE_INC },
            { FVIRTKEY | FALT, VK_LEFT, 100 },
            { FVIRTKEY | FALT, VK_RIGHT, 101 }
        };
        hAccel = CreateAcceleratorTable(acc, sizeof(acc) / sizeof(acc[0]));
    }
    
    void UpdateMenuChecks(HWND hwnd) {
        HMENU hBar = GetMenu(hwnd);
        if (!hBar) return;
        HMENU hSettings = GetSubMenu(hBar, 1);
        if (!hSettings) return;
        HMENU hTheme = GetSubMenu(hSettings, 0);
        HMENU hFont = GetSubMenu(hSettings, 1);
        HMENU hFontSize = GetSubMenu(hSettings, 2);
        
        ThemeChoice theme = Settings::GetTheme();
        UINT idTheme = (theme == ThemeChoice::System) ? ID_MENU_THEME_SYSTEM : 
                       (theme == ThemeChoice::Light ? ID_MENU_THEME_LIGHT : ID_MENU_THEME_DARK);
        if (hTheme) CheckMenuRadioItem(hTheme, ID_MENU_THEME_SYSTEM, ID_MENU_THEME_DARK, idTheme, MF_BYCOMMAND);
        
        std::wstring fontFace = Settings::GetFontFace();
        UINT idFont = ID_MENU_FONT_SEGOE;
        if (_wcsicmp(fontFace.c_str(), L"Consolas") == 0) idFont = ID_MENU_FONT_CONSOLAS;
        else if (_wcsicmp(fontFace.c_str(), L"Courier New") == 0) idFont = ID_MENU_FONT_COURIER;
        if (hFont) CheckMenuRadioItem(hFont, ID_MENU_FONT_SEGOE, ID_MENU_FONT_COURIER, idFont, MF_BYCOMMAND);
        
        int fontSize = Settings::GetFontSize();
        UINT idSize = (fontSize == 10) ? ID_MENU_FONTSIZE_10 : (fontSize == 12) ? ID_MENU_FONTSIZE_12 : 
                      (fontSize == 14) ? ID_MENU_FONTSIZE_14 : ID_MENU_FONTSIZE_16;
        if (hFontSize) CheckMenuRadioItem(hFontSize, ID_MENU_FONTSIZE_10, ID_MENU_FONTSIZE_16, idSize, MF_BYCOMMAND);
    }
}

// Navigation module

#ifndef UNICODE
#define UNICODE
#endif
#ifndef _UNICODE
#define _UNICODE
#endif

#include "../../include/ui.h"
#include "../../include/filesystem.h"
#include "../../include/settings.h"
#include <windows.h>
#include <commctrl.h>
#include <shellapi.h>
#include <algorithm>

namespace UI {
    void UpdateNavButtons(UIControls &controls, UIState &state) {
        EnableWindow(controls.hBackBtn, state.historyIndex > 0);
        EnableWindow(controls.hFwdBtn, state.historyIndex + 1 < (int)state.history.size());
    }
    
    void NavigateTo(HWND hwnd, UIControls &controls, UIState &state, 
                    const std::wstring &path, bool addToHistory) {
        
        if (addToHistory) {
            if (state.historyIndex + 1 < (int)state.history.size()) {
                state.history.erase(state.history.begin() + state.historyIndex + 1, state.history.end());
            }
            if (state.history.empty() || _wcsicmp(state.history.back().c_str(), path.c_str()) != 0) {
                state.history.push_back(path);
            }
            state.historyIndex = (int)state.history.size() - 1;
        }
        
        bool isNetworkPath = (path.length() >= 2 && path[0] == L'\\' && path[1] == L'\\');
        
        if (isNetworkPath) {
            HTREEITEM root = TreeView_GetRoot(controls.hTreeBottom);
            std::function<HTREEITEM(HTREEITEM, const std::wstring&)> searchTree = 
                [&](HTREEITEM hItem, const std::wstring& searchPath) -> HTREEITEM {
                    while (hItem) {
                        TVITEMW tvi{}; tvi.mask = TVIF_PARAM; tvi.hItem = hItem;
                        TreeView_GetItem(controls.hTreeBottom, &tvi);
                        Filesystem::NodeData *data = reinterpret_cast<Filesystem::NodeData*>(tvi.lParam);
                        if (data && _wcsicmp(data->path.c_str(), searchPath.c_str()) == 0) {
                            return hItem;
                        }
                        HTREEITEM child = TreeView_GetChild(controls.hTreeBottom, hItem);
                        if (child) {
                            HTREEITEM found = searchTree(child, searchPath);
                            if (found) return found;
                        }
                        hItem = TreeView_GetNextSibling(controls.hTreeBottom, hItem);
                    }
                    return nullptr;
                };
            HTREEITEM target = searchTree(root, path);
            if (target) {
                state.suppressTreeSelChange = true;
                TreeView_SelectItem(controls.hTreeBottom, target);
                TreeView_EnsureVisible(controls.hTreeBottom, target);
                state.suppressTreeSelChange = false;
            }
        } else {
            HTREEITEM target = Filesystem::EnsurePathInTree(controls.hTree, path);
            if (target) {
                state.suppressTreeSelChange = true;
                TreeView_SelectItem(controls.hTree, target);
                TreeView_EnsureVisible(controls.hTree, target);
                state.suppressTreeSelChange = false;
            }
        }
        
        state.currentPath = path;
        state.listItems.clear();
        
        // Update path edit IMMEDIATELY before loading directory to keep it synchronized
        if (controls.hPathEdit) {
            SetWindowTextW(controls.hPathEdit, path.c_str());
        }
        
        // Ensure list is in report mode with columns set up
        HWND hHeader = ListView_GetHeader(controls.hList);
        if (!IsReportMode(controls.hList)) {
            SetListViewMode(controls.hList, LVS_REPORT);
            DWORD ex = ListView_GetExtendedListViewStyle(controls.hList);
            ex |= LVS_EX_FULLROWSELECT | LVS_EX_LABELTIP;
            ThemeChoice effective = Settings::GetTheme();
            if (effective == ThemeChoice::System) {
                effective = Settings::QuerySystemAppsUseLightTheme() ? ThemeChoice::Light : ThemeChoice::Dark;
            }
            if (effective != ThemeChoice::Dark) {
                ex |= LVS_EX_GRIDLINES; // Only use system gridlines in light theme
            }
            ListView_SetExtendedListViewStyle(controls.hList, ex);
        } else {
            // Make sure columns exist - check header
            if (hHeader && Header_GetItemCount(hHeader) == 0) {
                SetListColumns(controls.hList);
            }
            AdjustReportColumns(controls.hList);
        }
        
        // Ensure image list is set before populating
        SHFILEINFOW sfi{};
        HIMAGELIST hSmall = (HIMAGELIST)SHGetFileInfoW(L"C:\\", FILE_ATTRIBUTE_DIRECTORY, &sfi, sizeof(sfi), SHGFI_SYSICONINDEX | SHGFI_SMALLICON);
        if (hSmall) {
            ListView_SetImageList(controls.hList, hSmall, LVSIL_SMALL);
        }
        
        Filesystem::PopulateListForDirectory(controls.hList, path, state.listItems, controls.hEmptyLabel);
        
        // Show/hide list view and header based on whether we have items
        if (!state.listItems.empty()) {
            ShowWindow(controls.hList, SW_SHOW);
            if (hHeader) ShowWindow(hHeader, SW_SHOW);
            if (controls.hEmptyLabel) {
                ShowWindow(controls.hEmptyLabel, SW_HIDE);
            }
        } else {
            // Hide list and header when showing empty message
            if (hHeader) ShowWindow(hHeader, SW_HIDE);
            ShowWindow(controls.hList, SW_HIDE);
        }
        
        UpdateNavButtons(controls, state);
        UpdateStatusBar(controls.hStatusBar, state.listItems);
        
        // Update WindowData with new state
        UI::UpdateWindowData(hwnd, controls, state);
    }
}

// Subclass procedures module

#ifndef UNICODE
#define UNICODE
#endif
#ifndef _UNICODE
#define _UNICODE
#endif

#include "../../include/ui.h"
#include "../../include/settings.h"
#include <windows.h>
#include <commctrl.h>

namespace UI {
    LRESULT CALLBACK ListViewSubclassProc(HWND hwnd, UINT uMsg, WPARAM wParam, LPARAM lParam) {
        HWND hParent = GetParent(hwnd);
        if (!hParent) return DefWindowProc(hwnd, uMsg, wParam, lParam);
        
        ThemeChoice effective = Settings::GetTheme();
        if (effective == ThemeChoice::System) {
            effective = Settings::QuerySystemAppsUseLightTheme() ? ThemeChoice::Light : ThemeChoice::Dark;
        }
        
        if (uMsg == WM_PAINT) {
            HWND hEmptyLabel = GetEmptyLabel(hParent);
            if (IsListViewEmpty(hParent) && hEmptyLabel && IsWindowVisible(hEmptyLabel)) {
                WNDPROC oldProc = GetOldListProc(hParent);
                if (!oldProc) return DefWindowProc(hwnd, uMsg, wParam, lParam);
                
                // Draw empty message
                LRESULT result = CallWindowProcW(oldProc, hwnd, uMsg, wParam, lParam);
                HDC hdc = GetDC(hwnd);
                if (hdc) {
                    RECT rc; GetClientRect(hwnd, &rc);
                    COLORREF grayColor = RGB(128, 128, 128);
                    if (effective == ThemeChoice::Dark) grayColor = RGB(160, 160, 160);
                    COLORREF oldColor = SetTextColor(hdc, grayColor);
                    int oldBkMode = SetBkMode(hdc, TRANSPARENT);
                    HFONT oldFont = nullptr;
                    HFONT hFont = Settings::GetFont();
                    if (hFont) oldFont = (HFONT)SelectObject(hdc, hFont);
                    wchar_t text[256];
                    GetWindowTextW(hEmptyLabel, text, 256);
                    if (wcslen(text) == 0) wcscpy(text, L"Здесь пусто");
                    SIZE textSize;
                    GetTextExtentPoint32W(hdc, text, (int)wcslen(text), &textSize);
                    int x = (rc.right - rc.left - textSize.cx) / 2;
                    int y = (rc.bottom - rc.top - textSize.cy) / 2;
                    TextOutW(hdc, x, y, text, (int)wcslen(text));
                    SetTextColor(hdc, oldColor);
                    SetBkMode(hdc, oldBkMode);
                    if (oldFont) SelectObject(hdc, oldFont);
                    ReleaseDC(hwnd, hdc);
                }
                return result;
            }
        }
        
        WNDPROC oldProc = GetOldListProc(hParent);
        if (!oldProc) return DefWindowProc(hwnd, uMsg, wParam, lParam);
        return CallWindowProcW(oldProc, hwnd, uMsg, wParam, lParam);
    }
    
    // Subclass procedure for border controls to handle dark theme borders
    LRESULT CALLBACK BorderSubclassProc(HWND hwnd, UINT uMsg, WPARAM wParam, LPARAM lParam, UINT_PTR uIdSubclass, DWORD_PTR dwRefData) {
        if (uMsg == WM_NCPAINT) {
            ThemeChoice effective = Settings::GetTheme();
            if (effective == ThemeChoice::System) {
                effective = Settings::QuerySystemAppsUseLightTheme() ? ThemeChoice::Light : ThemeChoice::Dark;
            }
            if (effective == ThemeChoice::Dark) {
                HDC hdc = GetWindowDC(hwnd);
                if (hdc) {
                    RECT rc;
                    GetWindowRect(hwnd, &rc);
                    OffsetRect(&rc, -rc.left, -rc.top);
                    // Draw dark border (gray)
                    COLORREF borderColor = RGB(80, 80, 80); // Gray border for dark theme
                    HBRUSH hBorderBrush = CreateSolidBrush(borderColor);
                    FrameRect(hdc, &rc, hBorderBrush);
                    DeleteObject(hBorderBrush);
                    ReleaseDC(hwnd, hdc);
                    return 0;
                }
            }
        }
        if (uMsg == WM_THEMECHANGED || uMsg == WM_SYSCOLORCHANGE) {
            // Invalidate to force redraw when theme changes
            InvalidateRect(hwnd, nullptr, FALSE);
        }
        return DefSubclassProc(hwnd, uMsg, wParam, lParam);
    }
    
    LRESULT CALLBACK HeaderSubclassProc(HWND hwnd, UINT uMsg, WPARAM wParam, LPARAM lParam) {
        HWND hParent = GetParent(hwnd);
        if (!hParent) return DefWindowProc(hwnd, uMsg, wParam, lParam);
        
        ThemeChoice effective = Settings::GetTheme();
        if (effective == ThemeChoice::System) {
            effective = Settings::QuerySystemAppsUseLightTheme() ? ThemeChoice::Light : ThemeChoice::Dark;
        }
        
        COLORREF bgColor = Settings::GetBackgroundColor();
        COLORREF textColor = Settings::GetTextColor();
        HBRUSH hBgBrush = CreateSolidBrush(bgColor);
        
        if (uMsg == WM_ERASEBKGND) {
            HDC hdc = (HDC)wParam;
            RECT rc; GetClientRect(hwnd, &rc);
            FillRect(hdc, &rc, hBgBrush);
            DeleteObject(hBgBrush);
            return 1;
        }
        if (uMsg == WM_PAINT) {
            PAINTSTRUCT ps;
            HDC hdc = BeginPaint(hwnd, &ps);
            if (hdc) {
                RECT rc; GetClientRect(hwnd, &rc);
                // Fill background
                FillRect(hdc, &rc, hBgBrush);
                
                if (effective == ThemeChoice::Dark) {
                    // Draw header items manually for dark theme
                    int itemCount = Header_GetItemCount(hwnd);
                    int x = 0;
                    for (int i = 0; i < itemCount; i++) {
                        HDITEMW hdi{};
                        wchar_t buf[256] = {0};
                        hdi.mask = HDI_TEXT | HDI_WIDTH | HDI_FORMAT;
                        hdi.pszText = buf;
                        hdi.cchTextMax = 256;
                        if (Header_GetItem(hwnd, i, &hdi)) {
                            RECT itemRect = {x, 0, x + hdi.cxy, rc.bottom};
                            // Draw separator line (gray)
                            if (i < itemCount - 1) {
                                HPEN hPen = CreatePen(PS_SOLID, 1, RGB(60, 60, 60));
                                HPEN hOldPen = (HPEN)SelectObject(hdc, hPen);
                                MoveToEx(hdc, itemRect.right - 1, itemRect.top, nullptr);
                                LineTo(hdc, itemRect.right - 1, itemRect.bottom);
                                SelectObject(hdc, hOldPen);
                                DeleteObject(hPen);
                            }
                            
                            // Draw text
                            SetTextColor(hdc, textColor);
                            SetBkColor(hdc, bgColor);
                            int oldBkMode = SetBkMode(hdc, TRANSPARENT);
                            DrawTextW(hdc, buf, -1, &itemRect, DT_LEFT | DT_VCENTER | DT_SINGLELINE);
                            SetBkMode(hdc, oldBkMode);
                            
                            x += hdi.cxy;
                        }
                    }
                } else {
                    // Use default drawing for light theme
                    WNDPROC oldProc = GetOldHeaderProc(hParent);
                    if (oldProc) {
                        CallWindowProcW(oldProc, hwnd, WM_PAINT, (WPARAM)hdc, 0);
                    }
                }
                
                EndPaint(hwnd, &ps);
            }
            DeleteObject(hBgBrush);
            return 0;
        }
        DeleteObject(hBgBrush);
        WNDPROC oldProc = GetOldHeaderProc(hParent);
        if (!oldProc) return DefWindowProc(hwnd, uMsg, wParam, lParam);
        return CallWindowProcW(oldProc, hwnd, uMsg, wParam, lParam);
    }
}

// Theme and font application module

#ifndef UNICODE
#define UNICODE
#endif
#ifndef _UNICODE
#define _UNICODE
#endif

#include "../../include/ui.h"
#include "../../include/settings.h"
#include <windows.h>
#include <commctrl.h>
#include <dwmapi.h>

// Forward declaration
extern LRESULT CALLBACK HeaderSubclassProc(HWND hwnd, UINT uMsg, WPARAM wParam, LPARAM lParam);

namespace UI {
    void ApplyThemeToControls(HWND hwnd, UIControls &controls) {
        COLORREF bgColor = Settings::GetBackgroundColor();
        COLORREF textColor = Settings::GetTextColor();
        
        if (controls.hTree) {
            TreeView_SetBkColor(controls.hTree, bgColor);
            TreeView_SetTextColor(controls.hTree, textColor);
            InvalidateRect(controls.hTree, nullptr, TRUE);
        }
        if (controls.hTreeBottom) {
            TreeView_SetBkColor(controls.hTreeBottom, bgColor);
            TreeView_SetTextColor(controls.hTreeBottom, textColor);
            InvalidateRect(controls.hTreeBottom, nullptr, TRUE);
        }
        if (controls.hSplitter) {
            InvalidateRect(controls.hSplitter, nullptr, TRUE);
        }
        if (controls.hList) {
            ThemeChoice effective = Settings::GetTheme();
            if (effective == ThemeChoice::System) {
                effective = Settings::QuerySystemAppsUseLightTheme() ? ThemeChoice::Light : ThemeChoice::Dark;
            }
            // For dark theme, use CLR_NONE to prevent background from overwriting items
            if (effective == ThemeChoice::Dark) {
                ListView_SetBkColor(controls.hList, CLR_NONE);
                ListView_SetTextBkColor(controls.hList, CLR_NONE);
            } else {
                ListView_SetBkColor(controls.hList, bgColor);
                ListView_SetTextBkColor(controls.hList, bgColor);
            }
            ListView_SetTextColor(controls.hList, textColor);
            if (effective == ThemeChoice::Dark) {
                HWND hHeader = ListView_GetHeader(controls.hList);
                if (hHeader) {
                    UI::SetHeaderProc(hwnd, hHeader, HeaderSubclassProc);
                    InvalidateRect(hHeader, nullptr, TRUE);
                }
            }
            InvalidateRect(controls.hList, nullptr, TRUE);
        }
        if (controls.hEmptyLabel) {
            InvalidateRect(controls.hEmptyLabel, nullptr, TRUE);
        }
        
        ThemeChoice effective = Settings::GetTheme();
        if (effective == ThemeChoice::System) {
            effective = Settings::QuerySystemAppsUseLightTheme() ? ThemeChoice::Light : ThemeChoice::Dark;
        }
        if (effective == ThemeChoice::Dark) {
            BOOL darkMode = TRUE;
            DwmSetWindowAttribute(hwnd, DWMWA_USE_IMMERSIVE_DARK_MODE, &darkMode, sizeof(darkMode));
            // Invalidate buttons for dark theme
            if (controls.hBackBtn) InvalidateRect(controls.hBackBtn, nullptr, TRUE);
            if (controls.hFwdBtn) InvalidateRect(controls.hFwdBtn, nullptr, TRUE);
            if (controls.hViewDetailsBtn) InvalidateRect(controls.hViewDetailsBtn, nullptr, TRUE);
            if (controls.hViewLargeBtn) InvalidateRect(controls.hViewLargeBtn, nullptr, TRUE);
            if (controls.hViewSmallBtn) InvalidateRect(controls.hViewSmallBtn, nullptr, TRUE);
            if (controls.hGoBtn) InvalidateRect(controls.hGoBtn, nullptr, TRUE);
            // Redraw menu bar
            DrawMenuBar(hwnd);
        } else {
            BOOL lightMode = FALSE;
            DwmSetWindowAttribute(hwnd, DWMWA_USE_IMMERSIVE_DARK_MODE, &lightMode, sizeof(lightMode));
            DrawMenuBar(hwnd);
        }
        if (controls.hPathEdit) {
            InvalidateRect(controls.hPathEdit, nullptr, TRUE);
        }
    }
    
    void ApplyFontToControls(HWND hwnd, UIControls &controls) {
        Settings::RecreateFont();
        HFONT hFont = Settings::GetFont();
        if (!hFont) return;
        
        if (controls.hTree) SendMessageW(controls.hTree, WM_SETFONT, (WPARAM)hFont, TRUE);
        if (controls.hTreeBottom) SendMessageW(controls.hTreeBottom, WM_SETFONT, (WPARAM)hFont, TRUE);
        if (controls.hList) {
            SendMessageW(controls.hList, WM_SETFONT, (WPARAM)hFont, TRUE);
            UpdateListViewRowHeight(controls.hList);
        }
        if (controls.hBackBtn) SendMessageW(controls.hBackBtn, WM_SETFONT, (WPARAM)hFont, TRUE);
        if (controls.hFwdBtn) SendMessageW(controls.hFwdBtn, WM_SETFONT, (WPARAM)hFont, TRUE);
        if (controls.hEmptyLabel) SendMessageW(controls.hEmptyLabel, WM_SETFONT, (WPARAM)hFont, TRUE);
        if (controls.hPathEdit) SendMessageW(controls.hPathEdit, WM_SETFONT, (WPARAM)hFont, TRUE);
        if (controls.hGoBtn) SendMessageW(controls.hGoBtn, WM_SETFONT, (WPARAM)hFont, TRUE);
        
        // Set font for menu bar
        HMENU hMenu = GetMenu(hwnd);
        if (hMenu && hFont) {
            // Menu font is handled by system, but we can force redraw
            NONCLIENTMETRICSW ncm = { sizeof(NONCLIENTMETRICSW) };
            if (SystemParametersInfoW(SPI_GETNONCLIENTMETRICS, sizeof(ncm), &ncm, 0)) {
                // Update system metrics would require SPI_SETNONCLIENTMETRICS which needs admin
                // Instead, force menu redraw
                DrawMenuBar(hwnd);
            }
        }
        
        // Relayout controls to match new font size
        UIState* state = UI::GetWindowStateForLayout(hwnd);
        if (state) {
            LayoutChildren(hwnd, controls, *state);
        }
    }
}

// Window data management module

#ifndef UNICODE
#define UNICODE
#endif
#ifndef _UNICODE
#define _UNICODE
#endif

#include "../../include/ui.h"
#include <windows.h>
#include <mutex>

// Window data structure to store UIControls and UIState
struct WindowData {
    UIControls controls;
    UIState state;
    std::mutex networkScanMutex;
};

// Store window data in GWLP_USERDATA
namespace {
    WindowData* GetWindowDataInternal(HWND hwnd) {
        return reinterpret_cast<WindowData*>(GetWindowLongPtrW(hwnd, GWLP_USERDATA));
    }
}

// Export functions for use in other modules
void SetWindowDataInternal(HWND hwnd, WindowData* data) {
    SetWindowLongPtrW(hwnd, GWLP_USERDATA, reinterpret_cast<LONG_PTR>(data));
}

// Export GetWindowData for use in other modules
WindowData* GetWindowData(HWND hwnd) {
    return GetWindowDataInternal(hwnd);
}

namespace UI {
    void SetWindowData(HWND hwnd, UIControls *controls, UIState *state) {
        WindowData* existing = GetWindowData(hwnd);
        if (existing) {
            if (controls) existing->controls = *controls;
            if (state) existing->state = *state;
        } else {
            WindowData* wndData = new WindowData;
            if (controls) wndData->controls = *controls;
            if (state) wndData->state = *state;
            SetWindowDataInternal(hwnd, wndData);
        }
    }
    
    UIControls* GetWindowControls(HWND hwnd) {
        WindowData* wndData = GetWindowData(hwnd);
        return wndData ? &wndData->controls : nullptr;
    }
    
    UIState* GetWindowState(HWND hwnd) {
        WindowData* wndData = GetWindowData(hwnd);
        return wndData ? &wndData->state : nullptr;
    }
    
    // Helper functions to update WindowData from external modules
    void UpdateWindowControls(HWND hwnd, const UIControls &controls) {
        WindowData* wndData = GetWindowData(hwnd);
        if (wndData) {
            wndData->controls = controls;
        }
    }
    
    void UpdateWindowState(HWND hwnd, const UIState &state) {
        WindowData* wndData = GetWindowData(hwnd);
        if (wndData) {
            wndData->state = state;
        }
    }
    
    void UpdateWindowData(HWND hwnd, const UIControls &controls, const UIState &state) {
        WindowData* wndData = GetWindowData(hwnd);
        if (wndData) {
            wndData->controls = controls;
            wndData->state = state;
        }
    }
    
    // Helper function to get state for layout
    UIState* GetWindowStateForLayout(HWND hwnd) {
        WindowData* wndData = GetWindowData(hwnd);
        return wndData ? &wndData->state : nullptr;
    }
    
    // Helper function to set header proc
    void SetHeaderProc(HWND hwnd, HWND hHeader, WNDPROC proc) {
        WindowData* wndData = GetWindowData(hwnd);
        if (wndData && hHeader && !wndData->state.oldHeaderProc) {
            wndData->state.oldHeaderProc = (WNDPROC)SetWindowLongPtrW(hHeader, GWLP_WNDPROC, (LONG_PTR)proc);
        }
    }
    
    // Helper functions for subclass procedures
    bool IsListViewEmpty(HWND hwnd) {
        WindowData* wndData = GetWindowData(hwnd);
        if (!wndData) return false;
        return wndData->state.listItems.empty();
    }
    
    HWND GetEmptyLabel(HWND hwnd) {
        WindowData* wndData = GetWindowData(hwnd);
        return wndData ? wndData->controls.hEmptyLabel : nullptr;
    }
    
    WNDPROC GetOldListProc(HWND hwnd) {
        WindowData* wndData = GetWindowData(hwnd);
        return wndData ? wndData->state.oldListProc : nullptr;
    }
    
    WNDPROC GetOldHeaderProc(HWND hwnd) {
        WindowData* wndData = GetWindowData(hwnd);
        return wndData ? wndData->state.oldHeaderProc : nullptr;
    }
    
    // Helper functions for control initialization
    WindowData* CreateWindowData() {
        return new WindowData;
    }
    
    void InitializeWindowData(WindowData* wndData, const UIControls &controls, const UIState &state) {
        if (wndData) {
            wndData->controls = controls;
            wndData->state = state;
        }
    }
    
    HWND GetToolTip(HWND hwnd) {
        WindowData* wndData = GetWindowData(hwnd);
        return wndData ? wndData->controls.hToolTip : nullptr;
    }
    
    void SetToolTip(HWND hwnd, HWND hToolTip) {
        WindowData* wndData = GetWindowData(hwnd);
        if (wndData) {
            wndData->controls.hToolTip = hToolTip;
        }
    }
    
    void UpdateScannedDevices(HWND hwnd, const std::vector<std::wstring> &devices) {
        WindowData* wndData = GetWindowData(hwnd);
        if (wndData) {
            std::lock_guard<std::mutex> lock(wndData->networkScanMutex);
            wndData->state.scannedDevices = devices;
        }
    }
    
    std::mutex* GetNetworkScanMutex(HWND hwnd) {
        WindowData* wndData = GetWindowData(hwnd);
        return wndData ? &wndData->networkScanMutex : nullptr;
    }
    
    std::vector<std::wstring>* GetScannedDevices(HWND hwnd) {
        WindowData* wndData = GetWindowData(hwnd);
        return wndData ? &wndData->state.scannedDevices : nullptr;
    }
    
    // Helper functions for network scan state
    HTREEITEM GetLoadingItem(HWND hwnd) {
        WindowData* wndData = GetWindowData(hwnd);
        return wndData ? wndData->state.loadingItem : nullptr;
    }
    
    void SetLoadingItem(HWND hwnd, HTREEITEM item) {
        WindowData* wndData = GetWindowData(hwnd);
        if (wndData) {
            wndData->state.loadingItem = item;
        }
    }
    
    bool GetNetworkScanInProgress(HWND hwnd) {
        WindowData* wndData = GetWindowData(hwnd);
        return wndData ? wndData->state.networkScanInProgress : false;
    }
    
    void SetNetworkScanInProgress(HWND hwnd, bool inProgress) {
        WindowData* wndData = GetWindowData(hwnd);
        if (wndData) {
            wndData->state.networkScanInProgress = inProgress;
        }
    }
    
    std::vector<std::wstring> GetScannedDevicesLocked(HWND hwnd) {
        WindowData* wndData = GetWindowData(hwnd);
        if (!wndData) return std::vector<std::wstring>();
        std::lock_guard<std::mutex> lock(wndData->networkScanMutex);
        return wndData->state.scannedDevices;
    }
    
    void ClearScannedDevices(HWND hwnd) {
        WindowData* wndData = GetWindowData(hwnd);
        if (wndData) {
            wndData->state.scannedDevices.clear();
        }
    }
    
    void DeleteWindowData(HWND hwnd) {
        WindowData* wndData = GetWindowData(hwnd);
        if (wndData) {
            delete wndData;
            SetWindowDataInternal(hwnd, nullptr);
        }
    }
}
\end{lstlisting}

\textbf{Utils}
\begin{lstlisting}[style=CodeListing]
#include "../../include/utils.h"
#include <shlwapi.h>
#include <algorithm>

namespace Utils {
    std::wstring JoinPath(const std::wstring &a, const std::wstring &b) {
        if (a.empty()) return b;
        if (a.back() == L'\\') return a + b;
        return a + L"\\" + b;
    }
    
    std::wstring FileTimeToWString(const FILETIME &ft) {
        SYSTEMTIME stUTC{}, stLocal{};
        FileTimeToSystemTime(&ft, &stUTC);
        SystemTimeToTzSpecificLocalTime(nullptr, &stUTC, &stLocal);
        wchar_t buf[64];
        swprintf(buf, 64, L"%04u-%02u-%02u %02u:%02u",
                 stLocal.wYear, stLocal.wMonth, stLocal.wDay, stLocal.wHour, stLocal.wMinute);
        return buf;
    }
}
\end{lstlisting}

\textbf{Main}
\begin{lstlisting}[style=CodeListing]
#ifdef _WIN32
#ifndef UNICODE
#define UNICODE
#endif
#ifndef _UNICODE
#define _UNICODE
#endif

// Winsock must be included before windows.h
#include <winsock2.h>
#include <ws2tcpip.h>
#include <windows.h>
#include <commctrl.h>
#include <shlwapi.h>
#include <gdiplus.h>
#include "include/ui.h"
#include "include/config.h"

#pragma comment(lib, "Comctl32.lib")
#pragma comment(lib, "gdiplus.lib")
#pragma comment(lib, "dwmapi.lib")
#pragma comment(lib, "mpr.lib")
#pragma comment(lib, "iphlpapi.lib")
#pragma comment(lib, "ws2_32.lib")
#pragma comment(lib, "shlwapi.lib")

using namespace Gdiplus;

// Initialize Winsock
struct WinsockInitializer {
    WinsockInitializer() {
        WSADATA wsaData;
        WSAStartup(MAKEWORD(2, 2), &wsaData);
    }
    ~WinsockInitializer() {
        WSACleanup();
    }
} g_winsockInit;

int APIENTRY WinMain(HINSTANCE hInstance, HINSTANCE, LPSTR, int nCmdShow) {
    // Initialize GDI+
    GdiplusStartupInput gdiplusStartupInput;
    ULONG_PTR gdiplusToken;
    GdiplusStartup(&gdiplusToken, &gdiplusStartupInput, nullptr);
    
    const wchar_t CLASS_NAME[] = L"FileManagerWin";
    WNDCLASSEXW wc{};
    wc.cbSize = sizeof(WNDCLASSEXW);
    wc.lpfnWndProc = UI::WindowProc;
    wc.hInstance = hInstance;
    wc.lpszClassName = CLASS_NAME;
    wc.hCursor = LoadCursor(nullptr, IDC_ARROW);
    wc.hbrBackground = (HBRUSH)(COLOR_WINDOW + 1);
    
    HICON hIcon = nullptr;
    HICON hIconSmall = nullptr;
    
    // Try to load icon from resource first
    hIcon = (HICON)LoadImageW(hInstance, MAKEINTRESOURCEW(IDI_MAINICON), IMAGE_ICON, 
                               GetSystemMetrics(SM_CXICON), GetSystemMetrics(SM_CYICON), 0);
    hIconSmall = (HICON)LoadImageW(hInstance, MAKEINTRESOURCEW(IDI_MAINICON), IMAGE_ICON,
                                    GetSystemMetrics(SM_CXSMICON), GetSystemMetrics(SM_CYSMICON), 0);
    
    if (!hIcon) {
        hIcon = LoadIconW(hInstance, MAKEINTRESOURCEW(IDI_MAINICON));
        hIconSmall = hIcon;
    }
    
    if (!hIcon) {
        // Try to load icon from icon.png using GDI+
        wchar_t exePath[MAX_PATH];
        GetModuleFileNameW(hInstance, exePath, MAX_PATH);
        PathRemoveFileSpecW(exePath);
        std::wstring iconPath = std::wstring(exePath) + L"\\icon.png";
        
        if (GetFileAttributesW(iconPath.c_str()) == INVALID_FILE_ATTRIBUTES) {
            iconPath = L"icon.png";
        }
        
        if (GetFileAttributesW(iconPath.c_str()) != INVALID_FILE_ATTRIBUTES) {
            hIcon = UI::LoadIconFromPNG(iconPath.c_str(), gdiplusToken);
            hIconSmall = hIcon;
        }
    }
    
    if (!hIcon) {
        hIcon = LoadIconW(nullptr, IDI_APPLICATION);
        hIconSmall = hIcon;
    }
    
    wc.hIcon = hIcon;
    wc.hIconSm = hIconSmall;
    
    if (!RegisterClassExW(&wc)) {
        GdiplusShutdown(gdiplusToken);
        return 1;
    }

    HWND hwnd = CreateWindowExW(0, CLASS_NAME, L"File Manager", WS_OVERLAPPEDWINDOW,
                                CW_USEDEFAULT, CW_USEDEFAULT, 1000, 700,
                                nullptr, nullptr, hInstance, nullptr);
    if (!hwnd) {
        GdiplusShutdown(gdiplusToken);
        return 1;
    }
    
    // Set window icon
    if (hIcon) {
        SendMessageW(hwnd, WM_SETICON, ICON_SMALL, (LPARAM)hIconSmall);
        SendMessageW(hwnd, WM_SETICON, ICON_BIG, (LPARAM)hIcon);
        SetClassLongPtrW(hwnd, GCLP_HICON, (LONG_PTR)hIcon);
        SetClassLongPtrW(hwnd, GCLP_HICONSM, (LONG_PTR)hIconSmall);
    }
    ShowWindow(hwnd, nCmdShow);
    UpdateWindow(hwnd);

    MSG msg{};
    while (GetMessage(&msg, nullptr, 0, 0)) {
        // Get accelerator table from window data (may not be available immediately)
        UIControls* controls = UI::GetWindowControls(hwnd);
        HACCEL hAccel = controls ? controls->hAccel : nullptr;
        
        if (!hAccel || !TranslateAccelerator(hwnd, hAccel, &msg)) {
            TranslateMessage(&msg);
            DispatchMessage(&msg);
        }
    }
    
    // Cleanup
    if (hIcon) DestroyIcon(hIcon);
    GdiplusShutdown(gdiplusToken);
    
    return (int)msg.wParam;
}

#else
// Linux/ncurses version - keep existing code
#include <ncurses.h>
// ... existing ncurses code ...
#endif
\end{lstlisting}